\chapter{预训练：从随机噪声到语言理解}
\label{ch:pretrain}
%% ============================================================

准备好了数据和模型架构，下一步就是预训练：
这是整个 LLM 训练流程中计算量最大、成本最高的阶段。
预训练的目标很简单：让模型学会预测下一个 token。
但实现这个「简单」目标所需的工程规模，
超乎大多数人的想象。

\section{训练目标：下一个 Token 预测}

LLM 的预训练使用\textbf{因果语言模型}（Causal Language Model）目标：
给定前面的所有 token，预测下一个 token。

形式化地，给定一个 token 序列 $x_1, x_2, \ldots, x_T$，
模型的训练损失是：
\begin{equation}
  \mathcal{L} = -\frac{1}{T} \sum_{t=1}^{T}
  \log P_\theta(x_t \mid x_1, x_2, \ldots, x_{t-1})
  \label{eq:lm-loss}
\end{equation}
这就是交叉熵损失，对数似然的负数，再对序列长度取平均。
模型被要求在每个位置输出一个词表大小的概率分布，
正确 token 的概率越高，损失越低。

\subsection{交叉熵、似然与 KL 散度的三位一体}

\cref{eq:lm-loss} 看起来是一个「工程选择」，
但它背后有深刻的信息论基础。
让我们逐步拆解这三个等价视角。

\textbf{第一视角：最大化对数似然。}
给定训练语料库中的所有句子，
我们希望找到参数 $\theta$ 使得模型「最可能生成」这些句子。
对于单个序列 $\mathbf{x} = (x_1, \ldots, x_T)$，
自回归分解给出的联合概率为：
\begin{equation}
  P_\theta(\mathbf{x})
  = \prod_{t=1}^{T} P_\theta(x_t \mid x_{<t})
  \label{eq:autoregressive-decomp}
\end{equation}
取对数得到 $\log P_\theta(\mathbf{x}) = \sum_{t=1}^T \log P_\theta(x_t \mid x_{<t})$。
最大化这个对数似然等价于最小化 \cref{eq:lm-loss}：
两者仅差一个负号和一个常数 $1/T$。

\textbf{第二视角：最小化 KL 散度。}
设 $P_{\mathrm{data}}$ 是训练数据的真实分布，
$P_\theta$ 是模型的分布。
我们可以写出 KL 散度：
\begin{equation}
  D_{\mathrm{KL}}(P_{\mathrm{data}} \| P_\theta)
  = \mathbb{E}_{P_{\mathrm{data}}} \left[
    \log \frac{P_{\mathrm{data}}(\mathbf{x})}{P_\theta(\mathbf{x})}
  \right]
  = H(P_{\mathrm{data}}, P_\theta) - H(P_{\mathrm{data}})
\end{equation}
其中 $H(P_{\mathrm{data}}, P_\theta)$ 是交叉熵，
$H(P_{\mathrm{data}})$ 是数据本身的熵（常数，不依赖 $\theta$）。
因此，\textbf{最小化交叉熵 $\Leftrightarrow$ 最小化 KL 散度
$\Leftrightarrow$ 最大化对数似然}。

这三条路殊途同归，但直觉各不相同：
对数似然视角说的是「让模型尽可能生成训练数据」；
KL 散度视角说的是「让模型分布尽可能接近真实分布」；
交叉熵视角说的是「让模型需要的平均编码长度最短」：
最终都指向同一个优化目标。

\textbf{第三视角：交叉熵的信息论解读。}
交叉熵 $H(P_{\mathrm{data}}, P_\theta)$ 可以理解为：
如果真实数据来自 $P_{\mathrm{data}}$，
但你使用模型 $P_\theta$ 来编码，
平均每个 token 需要多少比特。
好的语言模型 = 好的压缩器：
这正是 Shannon 信息论的核心洞察。
当模型完美时（$P_\theta = P_{\mathrm{data}}$），
交叉熵等于数据本身的熵，即理论上的最优编码长度。

\subsection{Teacher Forcing}

训练时，模型在预测第 $t$ 个 token 时
使用的是\textbf{真实的}前 $t-1$ 个 token 作为条件，
而不是模型自己之前生成的 token。
这种策略叫做 \textbf{Teacher Forcing}。

为什么不用模型自己的输出？
因为训练初期模型几乎是随机的，
如果用它自己的（错误的）预测作为后续的输入，
错误会指数级累积：
第 1 步预测错了，第 2 步基于错误输入预测更错，
第 3 步更加离谱。
这种「误差传播」会让训练完全无法收敛。

Teacher Forcing 通过「强制纠正」解决了这个问题：
不管模型第 $t-1$ 步预测得多差，
第 $t$ 步的输入始终是正确的 $x_{t-1}$。
这让每个位置的预测问题相互独立，
梯度信号清晰且有效。

代价是训练和推理之间存在分布偏移
（exposure bias）：
训练时模型总看到正确的前文，
推理时只能看到自己可能不正确的生成结果。
但在实践中，大规模预训练后这个问题的影响很小，
因为模型已经足够准确，
推理时的自回归生成很少偏离正轨。

\subsection{因果遮盖 vs 双向遮盖}

\textbf{因果遮盖}（Causal Masking）是 decoder-only 模型
（GPT 系列、LLaMA、DeepSeek 等）的核心机制：
在 Attention 计算中，
位置 $t$ 只能看到位置 $1$ 到 $t$ 的信息。
这通过一个下三角矩阵实现：
\begin{equation}
  \mathrm{Attn}(Q, K, V) = \mathrm{softmax}\left(
    \frac{QK^\top}{\sqrt{d}} + M
  \right) V, \quad
  M_{ij} = \begin{cases} 0 & i \geq j \\ -\infty & i < j \end{cases}
\end{equation}
$M$ 中的 $-\infty$ 在 softmax 后变成 0，
确保未来位置的信息不会泄漏到当前位置。

\textbf{双向遮盖}（Bidirectional Masking）则允许每个位置
看到序列中的所有其他位置：
这就是 BERT~\cite{devlin2019bert} 使用的方式。
双向模型在理解任务上有天然优势：
要判断一个句子的情感，当然应该同时看到前后文。

那为什么 decoder-only 最终赢了？
三个关键原因：
\begin{enumerate}[noitemsep]
  \item \textbf{生成能力}：因果遮盖天然支持自回归生成。
        双向模型无法逐步生成文本，因为它需要看到「未来」。
  \item \textbf{训练效率}：因果语言模型在每个位置都有预测目标，
        即每个 token 都贡献一个损失信号。
        BERT 的 MLM 只遮盖 15\% 的 token，
        85\% 的计算量没有直接贡献于训练目标。
  \item \textbf{Scaling 效果}：GPT-3~\cite{brown2020gpt3}
        和后续工作证明，在足够大的规模下，
        单向因果模型的理解能力并不逊于双向模型：
        通过预测下一个 token，模型被迫\textbf{隐式地}理解上下文。
\end{enumerate}

\begin{whybox}[为什么不用 BERT 的 MLM（遮盖语言模型）？]
BERT~\cite{devlin2019bert} 使用遮盖语言模型：
随机遮盖 15\% 的 token，让模型预测被遮盖的部分。
这是\textbf{双向}的，模型同时看到被遮盖位置的前后文。
这对理解任务（如问答、分类）很有帮助，
但有一个致命缺陷：它无法用于\textbf{生成}文本。
生成需要自回归过程——逐词生成，
每次只能看到前面已生成的内容。
因果语言模型天然支持这种自回归生成，
而且 GPT-3~\cite{brown2020gpt3} 证明了
单方向的因果模型在规模足够大时，
理解能力也不逊于双向模型。
因此，所有现代 LLM 都选择了因果语言模型。
\end{whybox}

\subsection{Prefix Language Modeling}

值得一提的是一种折中方案：
\textbf{Prefix Language Modeling}（前缀语言模型）。
它将输入序列分为两部分：
\textbf{前缀}（prefix）部分使用双向注意力（互相可见），
\textbf{生成}（generation）部分使用因果注意力（只看前文）。

Google 的 PaLM~\cite{chowdhery2022palm} 和 U-PaLM 都支持这种模式。
前缀部分通常是任务描述或上下文，
模型可以充分利用双向信息来理解它；
生成部分则自回归地产出答案。

这种设计试图兼得两者的优势：
前缀的双向理解 + 生成的自回归能力。
实现上，Prefix LM 使用一个修改过的注意力掩码：
前缀区域内的 token 使用全连接掩码（每个 token 可以看到所有其他 token），
而生成区域的 token 使用标准的因果掩码（只看前文）。
这与标准因果模型的唯一区别
在于注意力掩码矩阵的左上角块
从下三角变成了全 1 矩阵。
T5 和 UL2 也探索了类似的混合目标。

但在实践中，纯因果模型（decoder-only）的简洁性和 Scaling 效果
使它成为了事实标准：
额外的复杂性带来的收益不足以改变整个行业的选择。
前缀长度的选择增加了一个额外的超参数，
不同任务的最优前缀比例不同，
这使得统一的预训练变得复杂。
更关键的是，Raffel 等人（2020）在 T5 的大规模消融实验中
发现前缀 LM 的 scaling 行为不如纯因果 LM 稳定：
在小规模下前缀 LM 有优势，
但随着模型规模的增大，优势逐渐消失。
这一发现加速了行业向纯 decoder-only 架构的收敛。

\section{Scaling Laws：计算的经济学}
\label{sec:pretrain-scaling-laws}

\subsection{Kaplan Scaling Laws（2020）}

2020 年，OpenAI 发现了一个改变整个领域的定律~\cite{kaplan2020scaling}：
LLM 的性能（以损失值衡量）与模型大小、数据量和计算量
之间存在\textbf{精确的幂律关系}。

\begin{equation}
  L(N) \approx \left(\frac{N_c}{N}\right)^{\alpha_N}, \quad
  L(D) \approx \left(\frac{D_c}{D}\right)^{\alpha_D}, \quad
  L(C) \approx \left(\frac{C_c}{C}\right)^{\alpha_C}
  \label{eq:scaling}
\end{equation}
其中 $N$ 是模型参数量，$D$ 是训练数据量（token 数），
$C$ 是总计算量（FLOP），
$\alpha$ 和下标 $c$ 是经验常数。

Kaplan 等人拟合出的指数值为：
$\alpha_N \approx 0.076$，$\alpha_D \approx 0.095$，
$\alpha_C \approx 0.050$。

这些数字意味着什么？
以 $\alpha_N = 0.076$ 为例：
模型参数量每增加 \textbf{10 倍}，
损失值下降约 $10^{-0.076} \approx 0.84$，即下降 16\%。
这听起来不多，但在对数坐标下，
这条直线\textbf{没有拐点}：
不存在「增大模型就不管用了」的天花板。
从 100M 到 1B 到 10B 到 100B，
每增大 10 倍，损失都稳定下降约 16\%。

$\alpha_C = 0.050$ 的含义则更直接：
计算量（FLOP）每增加 10 倍，
损失下降约 $10^{-0.05} \approx 0.89$，即 11\%。
这给了整个行业巨大的信心：
只要你愿意投入足够的计算，就能造出更好的模型。

Kaplan 论文的另一个关键发现是：
当计算预算固定时，应该\textbf{优先增大模型、少用数据}。
他们认为参数量比数据量更重要：
$N$ 的增长速度应该快于 $D$。
这直接影响了 GPT-3 的设计选择：
175B 参数，但仅用 300B token 训练。

\begin{keybox}[Scaling Laws 的直觉类比]
想象你在装修一间房子：
\begin{itemize}[noitemsep]
  \item \textbf{模型参数量 $N$} $\approx$ 房子的面积，更大的房子能放更多东西。
  \item \textbf{训练数据量 $D$} $\approx$ 装修材料的数量，更多材料让房子更完善。
  \item \textbf{计算量 $C$} $\approx$ 装修工人的工时——预算总是有限的。
\end{itemize}
Kaplan 说：如果预算有限，先买大房子，材料少一点也行。
Chinchilla 说：不对，大房子配太少的材料是浪费：
房子面积和材料应该等比例增长。
\end{keybox}

\subsection{Chinchilla 定律（2022）}

2022 年，DeepMind 对 Scaling Laws 进行了重要修正~\cite{hoffmann2022chinchilla}。
Kaplan 等人认为应该优先增大模型、少用数据，
而 Chinchilla 论文证明：
\textbf{对于给定的计算预算，模型参数量和训练数据量应该等比例增长}。

具体而言，Chinchilla 定律建议每个参数大约需要 20 个 token 的训练数据。
这意味着一个 70B 参数的模型应该在约 1.4T token 上训练：
远多于 GPT-3（175B 参数但仅 300B token）所使用的数据量。
Chinchilla（70B 参数，1.4T token）在几乎所有基准测试上超越了 GPT-3。

Chinchilla 的核心发现可以更精确地表述为：
对于固定的计算预算 $C$，
最优的参数量 $N^*$ 和数据量 $D^*$ 满足：
\begin{equation}
  N^* \propto C^{0.50}, \quad D^* \propto C^{0.50}
  \label{eq:chinchilla-optimal}
\end{equation}
也就是说，计算预算翻倍时，
$N$ 和 $D$ 应该\textbf{各增长} $\sqrt{2} \approx 1.41$ 倍。
这与 Kaplan 的结论截然不同：
Kaplan 建议把大部分增量预算花在增大模型上，
Chinchilla 则说应该均分。

这个修正的影响是深远的。
按 Chinchilla 的标准，GPT-3 被严重「欠训练」了：
175B 参数的最优训练数据量应该是约 $175 \times 20 = 3.5$T token，
而 GPT-3 仅用了 300B token——差了超过 10 倍。
换句话说，如果把训练 GPT-3 的计算量换成训练一个
更小但更充分训练的模型（如 70B 配 1.4T token），
性能反而更好。

\subsection{后 Chinchilla 时代：过训练策略}
\label{sec:overtraining}

Chinchilla 定律回答的是「给定训练预算，如何最优地分配」。
但它有一个隐含假设：
它优化的是给定训练计算预算下的最优模型。
这忽略了一个现实：\textbf{训练是一次性的，推理是永久的}。

考虑两种方案：
\begin{itemize}[noitemsep]
  \item 方案 A：训练一个 70B 模型（Chinchilla 最优），1.4T token。
  \item 方案 B：训练一个 7B 模型，14T token（按 Chinchilla 过训练 10 倍）。
\end{itemize}
两者的训练计算量接近（$C \approx 6ND$），
但方案 B 的模型小 10 倍：
部署时每次推理的成本只有方案 A 的 $\sim$10\%。
如果这个模型要服务亿万用户，
方案 B 的总成本（训练 + 推理）可能远低于方案 A。

Meta 的 LLaMA~\cite{touvron2023llama} 系列开创了这一思路：
\begin{itemize}[noitemsep]
  \item LLaMA-1 7B：1T token（Chinchilla 建议约 140B token，过训练 $\sim$7$\times$）
  \item LLaMA 2 7B：2T token（过训练 $\sim$14$\times$）
  \item LLaMA 3 8B~\cite{dubey2024llama3}：15T token（过训练 $\sim$100$\times$）
\end{itemize}

DeepSeek-V3~\cite{deepseekv3} 正是基于这个思想：
671B 总参数（MoE 架构下每个 token 仅激活 37B），
在 14.8T token 上训练——远超 Chinchilla 最优数据量。
如果按 Chinchilla 的 20:1 比例，37B 活跃参数只需约 740B token，
实际用了 14.8T，过训练了约 20 倍。
这种「过训练」（over-training）策略使得模型在部署时极具成本优势。

\begin{corebox}[推理经济学：为什么「过训练」更划算]
假设训练成本为 $C_{\mathrm{train}}$，
模型每天推理 $Q$ 次查询，
每次推理成本正比于活跃参数量 $N_{\mathrm{active}}$，
模型服务周期为 $T$ 天。
总成本 =
$C_{\mathrm{train}} + Q \cdot T \cdot f(N_{\mathrm{active}})$。

对于广泛部署的模型（$Q \cdot T$ 极大），
推理项主导总成本：
此时减小 $N_{\mathrm{active}}$（即使需要增加 $C_{\mathrm{train}}$）
几乎总是有利的。
这就是 Chinchilla 最优不等于部署最优的根本原因。
\end{corebox}

\subsection{“Chinchilla is Dead”：推理感知的 Scaling Laws}
\label{sec:inference-aware-scaling}

2024--2025 年，学术界和工业界对 Chinchilla 定律展开了深入的反思，
甚至出现了 “Chinchilla is Dead” 的说法。
核心论点是：\textbf{Chinchilla 只优化了训练成本，
完全忽略了推理成本}：
而对于被广泛部署的模型，推理成本远超训练成本。

Sardana 等人~(2024) 在 ICML 上发表了% NEW REF: sardana2024inference = Beyond Chinchilla-Optimal: Accounting for Inference in Language Model Scaling Laws, Sardana et al., ICML 2024
形式化的\textbf{推理感知 Scaling Laws}。
他们将总成本建模为训练与推理的加权和：
\begin{equation}
  C_{\mathrm{total}} = C_{\mathrm{train}}(N, D)
  + R \cdot C_{\mathrm{inference}}(N)
  \label{eq:inference-aware}
\end{equation}
其中 $R$ 是模型生命周期内的总推理请求量。
当 $R$ 较大时（即模型被广泛部署），
最优策略是训练\textbf{更小的模型、用更多数据}：
这正是 LLaMA 系列和 DeepSeek 系列的做法。
Sardana 等人的分析表明，对于每天服务数百万查询的模型，
最优参数量可能只有 Chinchilla 最优的 $1/5$ 到 $1/10$，
但需要 $5\times$ 到 $50\times$ 的训练数据。

Bian 等人~(2025) 进一步将\textbf{模型架构}纳入 Scaling Laws 的考量。% NEW REF: bian2025inference = Scaling Inference-Efficient Language Models, Bian et al., arXiv 2025
他们发现，参数量相同的模型因架构差异
（如层数 vs 宽度、GQA vs MHA、FFN 比例等）
推理延迟最多相差 $3.5\times$。
这意味着 Scaling Laws 不应只优化参数量和数据量，
还应该\textbf{共同优化架构设计}。

\subsection{数据受限的 Scaling Laws 与合成数据}
\label{sec:data-constrained}

当高质量自然数据趋于耗尽时，
传统 Scaling Laws 的 “无限数据” 假设不再成立。

Muennighoff 等人~(2023) 的研究表明，% NEW REF: muennighoff2023scaling = Scaling Data-Constrained Language Models, Muennighoff et al., NeurIPS 2023
重复数据的价值递减但不为零：
数据重复 4 次的效果大约等价于使用 $4^{0.6} \approx 2.3$ 倍的去重数据。
超过约 4 个 epoch 后，重复的边际收益急剧下降。

这个问题催生了\textbf{合成数据 Scaling Laws} 的研究。
Qin 等人~(2025) 提出了 SynthLLM 框架，% NEW REF: qin2025synthllm = Scaling Laws of Synthetic Data for Language Models, Qin et al., arXiv 2025
首次系统地研究了合成数据的 scaling 行为。
他们发现：（1）朴素地扩大合成数据量存在严重的收益递减：
模型会过拟合于合成数据的分布偏差；
（2）通过“精心练习”（deliberate practice）策略，
即动态生成针对模型弱点的合成数据，
可以显著改善合成数据的 scaling 效率。
Microsoft 基于 SynthLLM 的实验表明，
经过精心设计的合成数据可以有效突破“数据墙”，
为模型提供接近真实数据质量的训练信号。

\subsection{Scaling Laws 的新方向（2025--2026）}

Scaling Laws 的研究仍在快速发展。
几个值得关注的方向：

\textbf{下游任务的 Scaling Laws}。
Held 等人~(2025) 提出了\textbf{相对 Scaling Laws}，% NEW REF: held2025relative = Relative Scaling Laws for LLMs, Held et al., arXiv 2025
不再用绝对的损失值来描述 scaling 行为，
而是预测模型之间的\textbf{相对性能差异}。
这种方法对不同评估基准的尺度变化更鲁棒。
Zhang 等人~(2026) 进一步提出了\textbf{规范性 Scaling Laws}，% NEW REF: zhang2026prescriptive = Prescriptive Scaling Reveals the Evolution of Language Model Capabilities, Zhang et al., arXiv 2026
基于 5000+ 个模型评估数据点，
估计给定预训练计算预算所能达到的下游准确率上界。
他们发现，随着后训练技术（RLHF、推理时计算等）的进步，
同样预训练预算对应的能力边界在持续上移：
这意味着 Scaling Laws 不是静态的，
它们会随着整个技术栈的进步而“演化”。

\textbf{推理时计算的 Scaling Laws}。
OpenAI 的 o1/o3 系列模型开辟了另一条 scaling 路径：
不增加模型参数，
而是在推理时投入更多计算（如更长的思维链、更多的采样和验证）。
这引出了一个新的权衡维度：
\textbf{训练时计算 vs 推理时计算}。
对于某些推理密集型任务（数学、编程），
在推理时使用更多计算的回报可能高于
同等计算量用于增大模型。

\subsection{Scaling Laws 的局限性}

尽管 Scaling Laws 是 LLM 领域最重要的实证发现之一，
它们有几个重要的局限性：

\textbf{不能预测涌现能力}（Emergent Capabilities）。
Scaling Laws 告诉你损失值会随规模平滑下降，
但不会告诉你模型在什么规模下突然学会了
思维链推理（Chain-of-Thought）、
上下文学习（In-Context Learning）或者代码生成。
这些能力似乎在某个规模阈值之后「突然出现」：
而这在平滑的幂律曲线上是看不到的。

\textbf{仅预测训练损失，不预测下游任务性能}。
训练损失（perplexity）的改善不一定线性地转化为
具体任务（如数学推理、代码生成）上的性能提升。
你可能需要将损失从 2.5 降到 2.3 才能看到数学能力的跳跃，
但从 2.3 降到 2.1 对数学帮助不大，
反而让诗歌写作能力有了飞跃。

\textbf{数据质量不在公式中}。
Scaling Laws 假设数据质量恒定，
但实际上数据质量可能比数据数量更重要。
Phi-2（2.7B）在精选的「教科书级」数据上训练，
在部分基准测试上超过了 25$\times$ 大的模型。

\textbf{「苦涩的教训」}。
Rich Sutton 的经典文章~\cite{sutton2019bitterlesson}
总结了 AI 研究 70 年的规律：
利用\textbf{更多计算}的方法，长期来看总是赢：
那些依赖人类知识和巧妙特征工程的方法，
短期内可能有效，但最终会被暴力扩展的方法超越。
Scaling Laws 是这个「苦涩教训」的数学表述。

\section{分布式训练：为什么需要数千块 GPU？}

让我们做一个简单的计算来理解为什么单块 GPU 远远不够。

一个 70B 参数的模型：
\begin{itemize}[noitemsep]
  \item \textbf{模型权重}：$70 \times 10^9 \times 2$ bytes (BF16) $= 140$\,GB
  \item \textbf{优化器状态}（Adam）：每个参数需要 12 bytes
        （FP32 参数副本 4 bytes + 一阶动量 4 bytes + 二阶动量 4 bytes）$= 840$\,GB
  \item \textbf{梯度}：$140$\,GB
  \item \textbf{激活值}：取决于 batch size 和序列长度，通常 100+\,GB
\end{itemize}
总计约 \textbf{1,220+\,GB}：
而一块 H100 GPU 只有 80\,GB 显存。
即使是最新的 H200（141\,GB），单块也放不下。

这就是为什么分布式训练是必须的。
现代 LLM 训练使用多种并行策略的组合。

\subsection{数据并行（Data Parallelism）}

最简单的并行方式：每块 GPU 持有完整的模型副本，
处理不同的数据批次，然后同步梯度。

\begin{lstlisting}[caption={Distributed Data Parallel in PyTorch}]
import torch.distributed as dist
from torch.nn.parallel import DistributedDataParallel as DDP

# Each GPU gets the same model
model = LLM().to(local_rank)
model = DDP(model, device_ids=[local_rank])

for batch in dataloader:
    loss = model(batch)
    loss.backward()
    # Gradients are automatically all-reduced across GPUs
    optimizer.step()
\end{lstlisting}

问题是：每块 GPU 都要放下完整模型。
对于 70B+ 的模型，这不可行。

\textbf{梯度同步}是数据并行的核心操作。
所有 GPU 需要在反向传播后将梯度「平均化」，
确保每块 GPU 上的模型更新一致。
最常用的算法是 \textbf{Ring AllReduce}：
$N$ 块 GPU 组成一个环，
每块 GPU 向下一块发送自己的梯度切片、
接收上一块的梯度切片，经过 $2(N-1)$ 步通信后，
所有 GPU 都拥有了完整的平均梯度。
Ring AllReduce 的通信量与 GPU 数量 $N$ 无关
（每块 GPU 发送和接收的数据量约为 $2 \cdot \frac{M}{N} \cdot (N-1) \approx 2M$，
其中 $M$ 是梯度总大小），
因此能很好地扩展到大规模集群。

\textbf{大 Batch 训练的挑战。}
数据并行天然地扩大了有效 batch size：
$N$ 块 GPU 每块处理 $B$ 个样本，有效 batch size 为 $NB$。
但大 batch 带来两个问题：
（1）学习率需要相应调整，否则收敛速度会下降：
通常使用线性缩放规则（linear scaling rule）：
$\mathrm{lr}_{\mathrm{effective}} = \mathrm{lr}_{\mathrm{base}} \times N$；
（2）过大的 batch 可能让模型陷入尖锐的极小值
（sharp minima），泛化性变差。
这就是为什么大 batch 训练几乎总是需要
\textbf{学习率预热}（warmup）：
让模型在初期用较小的步长探索参数空间。

\subsection{ZeRO 与 FSDP：分片数据并行}

DeepSpeed 的 ZeRO 优化~\cite{rajbhandari2020zero}
和 PyTorch 的 FSDP 解决了这个问题：
将模型的参数、梯度和优化器状态\textbf{分片}到不同 GPU 上，
每块 GPU 只持有 $1/N$ 的模型状态。
计算时，通过 All-Gather 临时收集需要的参数；
反向传播后，通过 Reduce-Scatter 分发梯度。

ZeRO 分为三个阶段，每个阶段分片的内容不同，
对应不同的内存节省和通信开销：

\begin{itemize}[noitemsep]
  \item \textbf{ZeRO-1}：只分片优化器状态。
  \item \textbf{ZeRO-2}：分片优化器状态 + 梯度。
  \item \textbf{ZeRO-3 / FSDP}：分片一切，包括模型参数。
\end{itemize}

让我们用一个 \textbf{10B 参数模型在 8 块 GPU 上}的例子
来量化每个阶段的内存节省：

\begin{center}\small
\renewcommand{\arraystretch}{1.3}
\begin{tabular}{lrrr}
  \toprule
  \rowcolor{figLightGray}
  \textbf{组件} & \textbf{无 ZeRO} & \textbf{每块 GPU 需存储} & \textbf{Bytes/param} \\
  \midrule
  模型参数（BF16） & 20\,GB & 每 GPU 全量 & 2 \\
  梯度（BF16） & 20\,GB & 每 GPU 全量 & 2 \\
  优化器状态（FP32） & & & \\
  \quad 参数副本 & 40\,GB & 每 GPU 全量 & 4 \\
  \quad 一阶动量 $m$ & 40\,GB & 每 GPU 全量 & 4 \\
  \quad 二阶动量 $v$ & 40\,GB & 每 GPU 全量 & 4 \\
  \midrule
  \textbf{总计} & \textbf{160\,GB} & & 16 \\
  \bottomrule
\end{tabular}
\end{center}

\textbf{ZeRO-1}：只分片优化器状态（120\,GB $\to$ 15\,GB/GPU），
每块 GPU 仍需 20+20+15 = \textbf{55\,GB}。
节省 $\sim$3$\times$。

\textbf{ZeRO-2}：分片优化器 + 梯度（120+20 = 140\,GB $\to$ 17.5\,GB/GPU），
每块 GPU 需 20+17.5 = \textbf{37.5\,GB}。
节省 $\sim$4$\times$。

\textbf{ZeRO-3}：分片一切（160\,GB $\to$ 20\,GB/GPU），
每块 GPU 需 \textbf{20\,GB}。
节省 $\sim$8$\times$（$= N$）。

PyTorch 的 FSDP（Fully Sharded Data Parallel）
本质上等价于 ZeRO Stage 3。
它将每个层的参数分片到所有 GPU，
前向传播时通过 All-Gather 临时收集完整参数，
计算完成后立即释放其他 GPU 的参数；
反向传播时同样 All-Gather 收集、计算梯度后
通过 Reduce-Scatter 将梯度分片回各 GPU。
这样每块 GPU 在任何时刻只持有一层的完整参数
加上 $1/N$ 的总模型状态，实现了极致的内存效率。

\subsubsection{FSDP vs DeepSpeed ZeRO：工程层面的差异}

虽然 FSDP 和 ZeRO-3 在\textbf{数学原理}上等价，
但工程实现上有重要差异，
直接影响实际训练的效率和易用性：

\begin{center}\small
\renewcommand{\arraystretch}{1.3}
\begin{tabular}{p{3cm}p{5.5cm}p{5.5cm}}
\toprule
\rowcolor{figLightGray}
\textbf{维度} & \textbf{PyTorch FSDP} & \textbf{DeepSpeed ZeRO-3} \\
\midrule
\textbf{参数分片粒度} &
  默认按“扁平化参数组”（FlatParameter）分片：
  将多个参数拼接成一个一维张量后均匀切分。
  FSDP2 改用逐参数分片，保留原始张量形状。 &
  按优化器组分片。
  支持更细粒度的分片策略控制。 \\
\textbf{通信 API} &
  基于 PyTorch DTensor 和 \texttt{c10d} 原生集合通信。
  与 PyTorch 生态（如 torch.compile）深度集成。 &
  自有通信后端。
  支持 NCCL 和自定义传输。
  与 PyTorch 核心耦合较松，
  但灵活性更高。 \\
\textbf{混合精度} &
  通过 PyTorch AMP 原生支持。
  FSDP 的 \texttt{MixedPrecision} 策略
  可以独立控制参数、梯度和通信的精度。 &
  通过 ZeRO++ 支持量化通信
  （INT8 通信 + FP16 计算），
  在跨节点通信中节省约 50\% 带宽。 \\
\textbf{检查点} &
  FSDP2 支持 DTensor 格式的检查点，
  可以跨不同并行配置加载
  （例如，4 GPU 保存的检查点可以在 8 GPU 上加载）。 &
  需要使用 DeepSpeed 的专有检查点格式。
  \texttt{universal\_checkpoint} 工具
  支持格式转换，但增加了工程复杂度。 \\
\textbf{与 TP/PP 组合} &
  需要与 Megatron-LM 或
  其他 TP/PP 框架手动集成。
  Megatron-FSDP（2025）简化了这一过程。 &
  DeepSpeed 3D 原生支持
  DP + TP + PP 的组合。
  但 TP 实现的优化程度不如 Megatron-LM。 \\
\bottomrule
\end{tabular}
\end{center}

\textbf{FSDP2}（PyTorch 2.4+）解决了 FSDP 初代版本的
多个工程痛点。
最关键的改进是\textbf{逐参数分片}：
不再将参数扁平化为一维张量，
而是保留每个参数的原始形状和元信息。
这对\textbf{矩阵感知优化器}（如 Muon、Shampoo）至关重要，
因为这些优化器需要知道参数的二维矩阵结构
来执行正交化或 Kronecker 积运算：
如果参数被扁平化为一维向量，这些信息就丢失了。
veScale-FSDP 更进一步，
为非标准优化器和块级量化提供了“结构感知分片”：
这是 Kimi K2 等使用 Muon 优化器的模型
能够在分布式环境中高效训练的关键。

\textbf{ZeRO++ }（DeepSpeed，2023）
在 ZeRO-3 基础上增加了三项优化：
(1) \textbf{量化权重通信}：
All-Gather 时使用 INT8 而非 BF16 传输参数，
节省 50\% 的通信带宽，
到达后解量化回 BF16；
(2) \textbf{分层分区}：
在节点内保留额外的参数副本，
减少跨节点通信的频率；
(3) \textbf{量化梯度通信}：
Reduce-Scatter 使用 INT4/INT8 梯度。
这三项优化组合后，
在跨节点带宽受限的环境中
（如 DeepSeek 使用的 H800 集群，
跨节点带宽远低于无限制 H100 集群），
ZeRO++ 可以将端到端训练速度提升 \textbf{2.2$\times$}。

\textbf{选择建议（2026 年）}：
对于从零开始的新项目，
FSDP2 + Megatron-FSDP 是推荐方案：
它与 PyTorch 生态的集成最紧密，
检查点管理最灵活，
且对新优化器（Muon）的支持最好。
对于需要在带宽受限环境下运行的项目
（特别是使用受出口管制限制的 GPU 时），
DeepSpeed ZeRO++ 的量化通信优势不可替代。
对于 MoE 架构，
DeepSpeed 的专家并行支持更成熟，
但 Megatron Core MoE 正在迅速追赶。

\subsection{张量并行（Tensor Parallelism）}

将单个矩阵运算拆分到多块 GPU 上。
例如，一个 $[d, 4d]$ 的 FFN 权重矩阵可以沿列切分到 4 块 GPU，
每块计算 $[d, d]$ 的子矩阵乘法，
然后通过 All-Reduce 合并结果。

Megatron-LM 定义了两种切分方式：

\textbf{列并行}（Column Parallel）：
将权重矩阵 $W \in \mathbb{R}^{d \times kd}$ 沿列切分为
$W = [W_1, W_2, \ldots, W_k]$，
每块 GPU $i$ 计算 $Y_i = XW_i$。
由于激活函数（如 GeLU）是逐元素操作，
可以在切分后独立应用，不需要通信。
最终通过 All-Reduce 合并结果。

\textbf{行并行}（Row Parallel）：
将权重矩阵沿行切分为 $W = [W_1; W_2; \ldots; W_k]^T$，
输入也对应切分。
每块 GPU 计算部分结果后通过 All-Reduce 求和。

Megatron-LM 的巧妙之处在于
将 FFN 层的第一个线性层做列并行、
第二个线性层做行并行：
这样两层之间不需要额外通信，
只在最终输出时做一次 All-Reduce。

让我们通过一个具体例子来理解这种切分策略。
考虑一个 FFN 层：$Y = \mathrm{GeLU}(XW_1) W_2$，
其中 $X \in \mathbb{R}^{B \times d}$，
$W_1 \in \mathbb{R}^{d \times 4d}$，
$W_2 \in \mathbb{R}^{4d \times d}$。
使用 TP=4（4 块 GPU）：

\textbf{第一个线性层（列并行）}：
$W_1$ 沿列切分为 $[W_1^{(1)}, W_1^{(2)}, W_1^{(3)}, W_1^{(4)}]$，
每块 $W_1^{(i)} \in \mathbb{R}^{d \times d}$。
GPU $i$ 计算 $Y_i = \mathrm{GeLU}(X W_1^{(i)})$。
关键点：GeLU 是逐元素操作，
可以在分片后独立应用，\textbf{不需要通信}。
这是列并行的核心优势：
如果激活函数是 “通道内” 操作
（只依赖单个通道的值），
列切分后就可以独立计算。

\textbf{第二个线性层（行并行）}：
$W_2$ 沿行切分为 $[W_2^{(1)}; W_2^{(2)}; W_2^{(3)}; W_2^{(4)}]^\top$，
每块 $W_2^{(i)} \in \mathbb{R}^{d \times d}$。
GPU $i$ 计算部分结果 $Z_i = Y_i W_2^{(i)}$。
最终结果需要\textbf{求和}：$Z = \sum_i Z_i$。
这个求和通过 All-Reduce 完成：
这是整个 FFN 层中\textbf{唯一一次通信}。

因此，一个完整的 FFN 层在 TP=4 下
只需要\textbf{2 次 All-Reduce}
（前向和反向各 1 次），
而非朴素实现中的 4 次。
这种“列-行交替”的切分策略
是 Megatron-LM 效率的核心所在。

\textbf{Attention 头的切分}也是张量并行的一部分。
Multi-Head Attention 天然地将头分配到不同 GPU：
如果有 32 个头、4 块 GPU，
每块 GPU 处理 8 个头的 $Q$、$K$、$V$ 投影和 Attention 计算，
最终通过 All-Reduce 合并输出。
GQA（Grouped-Query Attention）对 TP 友好度更高：
因为 KV 头的数量较少（如 8 个），
可以在 TP 组内共享而不增加通信。

\textbf{序列并行}（Sequence Parallelism，SP）
是张量并行的一个重要补充。
在 TP 切分的层中，
LayerNorm 和 Dropout 等\textbf{逐元素操作}
不受 TP 的影响：
它们在每块 GPU 上独立处理完整的隐状态。
但这意味着每块 GPU 上这些操作的\textbf{激活值}
是完整的（未分片的），
在长序列场景下占据大量显存。

序列并行将这些“非 TP 操作”的激活值
沿\textbf{序列维度}分片到 TP 组的各 GPU 上。
具体实现是：在 TP 层的 All-Reduce 之前，
先做一次 Reduce-Scatter（结果被分片到各 GPU），
每块 GPU 只对自己负责的序列片段
执行 LayerNorm 和 Dropout；
在进入下一个 TP 层时，
通过 All-Gather 恢复完整的序列。
这使得 LayerNorm 和 Dropout 的激活值
也被分片了，显存节省约 $1/\mathrm{TP}$。

张量并行要求 GPU 之间有极高带宽的通信
（每次前向和反向传播都需要同步），
因此通常只在同一节点内的 GPU 之间使用（通过 NVLink 连接）。
跨节点的 InfiniBand（即使是 800\,Gb/s）也远不如
节点内 NVLink（H100 提供 900\,GB/s 双向带宽）。
这就是为什么张量并行的度数通常等于
单节点内的 GPU 数量（DGX H100 为 8）。

\textbf{TP 的通信开销}可以精确计算。
对于一个 Transformer 层，
TP 引入的通信量为每个 TP 层 2 次 All-Reduce，
每次 All-Reduce 传输 $2 \times B \times S \times d$ 字节
（$B$ = batch size, $S$ = 序列长度, $d$ = 隐藏维度，因子 2 来自 Ring AllReduce）。
一个典型的 70B 模型（$d = 8192$, $S = 4096$）
在 TP=8 下，每步每层的 TP 通信量约为
$4 \times B \times 4096 \times 8192 \times 2$ bytes (BF16)
$= B \times 268$\,MB。
在 NVLink 900\,GB/s 的带宽下，
这需要约 $B \times 0.3$\,ms：
对于 $B = 1$ 这几乎可以忽略，
但当 $B$ 较大或序列更长时，
TP 通信会成为\textbf{非平凡的瓶颈}。

\subsection{流水线并行（Pipeline Parallelism）}

将模型的不同层分配到不同 GPU 上，
数据像流水线一样依次通过各 GPU。
GPU 1 处理第 1--10 层，GPU 2 处理第 11--20 层，以此类推。

\textbf{流水线气泡问题}是核心挑战：
如果不做特殊处理，当 GPU 1 在处理前向传播时，
GPU 2--4 在等待 GPU 1 的输出；
当 GPU 4 在做反向传播时，GPU 1--3 在等待。
在简单的流水线调度中，气泡率为 $(P-1)/M$，
其中 $P$ 是流水线阶段数，$M$ 是 micro-batch 数。

\textbf{GPipe}~通过将一个 mini-batch 切分为 $M$ 个 micro-batch
来减少气泡：多个 micro-batch 依次进入流水线，
让各 GPU 尽可能保持忙碌。
但 GPipe 仍然需要先完成所有前向再开始所有反向。% NEW REF: gpipe = GPipe: Efficient Training of Giant Neural Networks using Pipeline Parallelism, Huang et al., NeurIPS 2019

\textbf{1F1B}（One Forward, One Backward）调度更进一步：
每完成一个 micro-batch 的前向传播，
立即开始另一个 micro-batch 的反向传播。
这种交错调度减少了内存峰值
（不需要同时保存所有 micro-batch 的激活值），
同时保持了流水线的利用率。

\textbf{Megatron-LM 的交错调度}（Interleaved Schedule）
进一步减少气泡：每个 GPU 不只负责连续的层，
而是负责多个不连续的层段（virtual pipeline stages）。
例如，4 块 GPU、32 层的模型，
传统方式是 GPU 1 = 层 1--8，GPU 2 = 层 9--16，...；
交错方式是 GPU 1 = 层 1--4 和 17--20，
GPU 2 = 层 5--8 和 21--24，以此类推。
这使得前向和反向的路径更短，气泡更小。

DeepSeek 的 DualPipe 技术进一步优化了流水线调度：
将前向和反向计算交错排列，
使得一个 micro-batch 的前向与另一个 micro-batch 的反向
可以在不同 GPU 上并行进行，
将流水线气泡率降低到约 3\%（传统方法通常 20\%+）。

\subsection{专家并行（Expert Parallelism）}

对于 MoE 模型（详见\cref{ch:moe}），
不同的专家分布在不同的 GPU 上。
当一个 token 被路由到某个专家时，
它需要通过 All-to-All 通信发送到对应的 GPU，
计算完成后再发送回来。

\textbf{All-to-All 通信}是专家并行独有的通信模式。
与数据并行的 AllReduce（每块 GPU 发送和接收相同大小的数据）不同，
All-to-All 的通信量取决于路由决策：
如果某个专家特别「热门」，
持有该专家的 GPU 会收到不成比例的大量 token。
这就是为什么 MoE 的负载均衡（\cref{ch:moe}）如此重要。

\textbf{容量因子}（Capacity Factor）是一个实用的工程参数：
它规定每个专家最多能处理 $\mathrm{CF} \times \frac{T}{E}$ 个 token
（$T$ = 总 token 数，$E$ = 专家数，$\mathrm{CF}$ 通常取 1.0--1.5）。
超出容量的 token 会被\textbf{丢弃}（token dropping），
直接跳过 MoE 层，使用残差连接的输出。
训练时适度的 token dropping 对质量影响不大，
但它保证了计算和通信的负载是有界的。
DeepSeek 通过无辅助损失的负载均衡
大幅减少了 token dropping 的发生。

DeepSeek 开发的 DeepEP 库专门优化了这种通信模式。

\subsection{上下文并行（Context Parallelism）}

当序列长度超过单块 GPU 能处理的范围时
（例如 128K 或更长的上下文），
需要将序列本身拆分到多块 GPU 上：
这就是上下文并行。

\textbf{Ring Attention} 是最著名的方案：
将长序列的 $Q$、$K$、$V$ 切分到 $P$ 块 GPU 上，
每块 GPU 持有序列的 $1/P$ 段。
$K$ 和 $V$ 在 GPU 之间以环形方式传递：
每个 step 中，GPU $i$ 计算自己的 $Q$ 与当前 $K$、$V$ 的注意力，
同时将 $K$、$V$ 发送给下一块 GPU。
经过 $P$ 步后，每块 GPU 都计算了与所有位置的注意力。
通信可以与计算完全重叠，效率很高。

\textbf{Ulysses}（DeepSpeed-Ulysses）采用了不同的切分策略：
它沿着\textbf{注意力头}的维度切分序列：
每块 GPU 处理所有位置但只处理部分头。
前向传播时通过 All-to-All 通信交换数据，
使得每块 GPU 获得完整的头信息。
Ulysses 对 Grouped-Query Attention（GQA）友好，
通信量更小，适合中等长度的序列。

\subsection{多维并行组合}
\label{sec:5d-parallelism}

\cref{fig:3d-parallelism} 展示了经典的 3D 并行
（DP $\times$ TP $\times$ PP）如何将训练工作
分配到一个 GPU 的三维网格上。
理解这个图是理解大规模训练工程的关键：
每一维对应一种并行策略，
每一维对通信带宽和延迟有不同的要求。

\begin{figure}[htbp]
\centering
\begin{tikzpicture}[
  % Isometric axes: x goes right, y goes up-right (depth), z goes up
  x={(1.0cm, 0cm)},
  y={(0.5cm, 0.38cm)},
  z={(0cm, 1.0cm)},
  gpu/.style={
    rectangle, rounded corners=1pt,
    minimum width=9mm, minimum height=9mm,
    inner sep=1pt, line width=0.6pt,
    font=\scriptsize\ttfamily
  },
]

% --- Bottom layer (TP rank 0) ---
\foreach \dp in {0,...,3} {
  \foreach \pp in {0,...,3} {
    \node[gpu, draw=figBlue!50, fill=figBlue!8, text=figBlue!80!black]
      (b-\dp-\pp) at (\dp*2.0, \pp*1.6, 0)
      {\dp,\pp,0};
  }
}

% --- Top layer (TP rank 1) ---
\foreach \dp in {0,...,3} {
  \foreach \pp in {0,...,3} {
    \node[gpu, draw=figRed!50, fill=figRed!8, text=figRed!80!black]
      (t-\dp-\pp) at (\dp*2.0, \pp*1.6, 1.8)
      {\dp,\pp,1};
  }
}

% --- TP connections (vertical, green dashed) ---
\foreach \dp in {0,...,3} {
  \foreach \pp in {0,...,3} {
    \draw[figGreen!50, line width=0.5pt, densely dashed]
      (b-\dp-\pp.north) -- (t-\dp-\pp.south);
  }
}

% --- PP connections (depth direction, orange arrows) ---
% Bottom layer only (top layer mirrors)
\foreach \dp in {0,...,3} {
  \foreach \pp in {0,...,2} {
    \pgfmathtruncatemacro{\ppnext}{\pp+1}
    \draw[figOrange!60, line width=0.5pt,
          -{Stealth[length=2pt,width=2pt]}]
      (b-\dp-\pp.north east) -- (b-\dp-\ppnext.south west);
  }
}

% --- DP indication: highlight one row with blue line ---
\draw[figBlue!40, line width=0.5pt, densely dotted]
  (b-0-0.west) -- (b-1-0.west) -- (b-2-0.west) -- (b-3-0.east);

% --- Axis arrows ---
\draw[figGray!50, line width=0.6pt, -{Stealth[length=2.5pt,width=2.5pt]}]
  (-0.6, -0.5, 0) -- (8.2, -0.5, 0)
  node[pos=1.0, below left, font=\scriptsize, text=figBlue!80!black]
  {DP $\times 4$};
\draw[figGray!50, line width=0.6pt, -{Stealth[length=2.5pt,width=2.5pt]}]
  (-0.6, -0.5, 0) -- (-0.6, 5.8, 0)
  node[pos=1.0, above left, font=\scriptsize, text=figOrange!80!black, rotate=37]
  {PP $\times 4$};
\draw[figGray!50, line width=0.6pt, -{Stealth[length=2.5pt,width=2.5pt]}]
  (-0.6, -0.5, 0) -- (-0.6, -0.5, 2.6)
  node[pos=1.0, left, font=\scriptsize, text=figGreen!80!black]
  {TP $\times 2$};

% --- Legend box ---
\node[anchor=north west, draw=figGray!25, rounded corners=2pt,
      fill=white, inner sep=5pt, line width=0.5pt,
      font=\scriptsize, align=left]
  at (8.8, 5.5, 0) {%
    \textcolor{figGreen!80!black}{\rule[0.5ex]{6pt}{0.6pt}\kern1pt\rule[0.5ex]{2pt}{0.6pt}\kern1pt\rule[0.5ex]{6pt}{0.6pt}}
    ~TP: NVLink (900\,GB/s)\\[2pt]
    \textcolor{figOrange!80!black}{$\longrightarrow$}
    ~PP: NVLink/IB (100--900)\\[2pt]
    \textcolor{figBlue!80!black}{$\cdots\cdots$}
    ~DP: AllReduce (IB, 100)
  };

\end{tikzpicture}
\caption{3D parallelism layout: $\mathrm{DP} \times \mathrm{PP} \times \mathrm{TP} = 4 \times 4 \times 2 = 32$ GPUs.
Each cell is one GPU labeled $(d, p, t)$.
\textcolor{figGreen}{Green dashed lines}: TP communication (NVLink, highest bandwidth).
\textcolor{figOrange}{Orange arrows}: PP communication (activations/gradients across pipeline stages).
DP AllReduce occurs along the horizontal axis among identically-positioned GPUs.
Bandwidth requirement: TP $>$ PP $>$ DP.}
\label{fig:3d-parallelism}
\end{figure}

实际训练中，这些并行策略是组合使用的。
DeepSeek-V3 的训练使用了\textbf{多维并行}：
\begin{enumerate}[noitemsep]
  \item \textbf{DP}（数据并行）：不同 GPU 组处理不同数据。
  \item \textbf{TP}（张量并行）：节点内 GPU 拆分矩阵运算。
  \item \textbf{PP}（流水线并行）：不同节点处理不同层。
  \item \textbf{SP}（序列并行）：拆分长序列的归一化和 Dropout 操作
        （SP 通常与 TP 配合使用，
        在 TP 切分的基础上沿序列维度分摊非张量并行操作的内存，
        因此有时不被视为独立的并行维度）。
  \item \textbf{EP}（专家并行）：MoE 的专家分布在不同 GPU 上。
\end{enumerate}

在 2048 块 H800 GPU 上协调这五个维度，是一项极其复杂的工程挑战。

让我们具体看一个例子来理解这些并行维度如何组合。
假设我们有 2048 块 GPU，训练一个 671B MoE 模型
（如 DeepSeek-V3），一种可能的配置是：
\begin{itemize}[noitemsep]
  \item \textbf{TP = 4}：每 4 块 GPU 协同处理一个层内的矩阵运算
        （通过 NVLink 连接，带宽最高）。
  \item \textbf{PP = 8}：模型的不同层分布在 8 个流水线阶段
        （跨节点通信，使用 InfiniBand）。
  \item \textbf{EP = 4}：256 个 MoE 专家分布在 4 个专家并行组
        （需要 All-to-All 通信）。
  \item \textbf{DP = 16}：剩余的 $2048 / (4 \times 8 \times 4) = 16$ 份
        用于数据并行（最简单的通信模式）。
\end{itemize}

这意味着系统被组织成一个四维网格：
$16 \times 8 \times 4 \times 4 = 2048$。
每一个维度对通信带宽和延迟有不同的要求，
这决定了它们在物理拓扑中的位置：
TP 在节点内（NVLink），
PP 在节点间但延迟敏感，
EP 需要高带宽但可以容忍一定延迟，
DP 对带宽和延迟最宽容。

另一个参考点是 \textbf{LLaMA 3 405B}~\cite{dubey2024llama3}
在 16384 块 H100 GPU 上的配置，使用了 \textbf{4D 并行}：
TP = 8，CP = 16，PP = 16，DP = 8
（$8 \times 16 \times 16 \times 8 = 16{,}384$）。
这里 TP = 8 对应 DGX H100 节点的 8 块 GPU
（全部通过 NVLink 连接），
CP = 16（Context Parallelism）用于支持 128K 长上下文训练，
将序列拆分到 16 块 GPU 上以分摊 KV cache 内存，
PP = 16 将 126 层分为 16 个阶段（每阶段约 8 层），
DP = 8 提供数据并行度以维持足够的全局 batch size
（约 16M token per batch）。

\subsection{通信--计算重叠}

分布式训练的效率不仅取决于「如何拆分计算」，
更取决于「如何隐藏通信开销」。
理想状态是：当 GPU A 在计算时，
GPU B 在同时发送上一步的结果：
通信完全被计算「遮盖」，
总时间等于纯计算时间。

实现这一点需要精细的调度。
以 ZeRO-3 / FSDP 为例：
当 GPU 在计算第 $i$ 层的前向传播时，
它可以提前发起第 $i+1$ 层参数的 All-Gather 请求。
等第 $i$ 层计算完毕，第 $i+1$ 层的参数已经准备好了，
GPU 可以无缝开始第 $i+1$ 层的计算。
反向传播时同理：完成第 $i$ 层的梯度后，
立即发起 Reduce-Scatter，
同时开始第 $i-1$ 层的梯度计算。

\textbf{梯度分桶}（Gradient Bucketing）是另一个关键优化。
PyTorch DDP 不会等所有梯度计算完毕后再发起 AllReduce：
它将参数分成多个「桶」（bucket），
每个桶的梯度计算完毕后立即开始 AllReduce。
反向传播是从最后一层开始的，
所以最后一层的梯度最先准备好：
当中间层还在做反向计算时，
最后几层的梯度已经在通信了。
好的分桶策略可以让通信和计算几乎完全重叠。

DeepEP 的 Low-Latency 模式更进一步：
它使用 GPU 硬件的 hook 机制
将通信操作从 GPU 的 SM（Streaming Multiprocessor）中完全卸载，
通信不占用任何计算资源，
实现了真正的零开销通信。

DeepSeek 的 \textbf{DualPipe} 将通信--计算重叠
推广到了流水线并行的维度。
传统的 1F1B 调度中，
前向的通信（接收上一阶段的激活值）和计算是串行的；
DualPipe 将流水线的前向和反向
拆分为更细粒度的「通信阶段」和「计算阶段」，
让来自不同 micro-batch 的通信和计算在时间上交错。
其核心思想是：当 GPU 在为 micro-batch $A$ 做计算时，
同时为 micro-batch $B$ 做通信：
因为计算密集用的是 SM，通信密集用的是 NVLink/IB 控制器，
两者可以并行而不冲突。
这让 DeepSeek-V3 在 2048 块 H800 上
实现了仅 $\sim$3\% 的流水线气泡率。

\subsection{训练框架演进（2025--2026）}

分布式训练框架在 2025--2026 年经历了重要的融合与进化。

\textbf{Megatron-FSDP}
是 NVIDIA 在 2025 年推出的新一代分布式训练方案，% NEW REF: megatronfsdp = Megatron-FSDP: Performance-oriented FSDP for Megatron-LM, NVIDIA, 2025
将 Megatron-LM 的高性能张量并行
与 PyTorch FSDP 的灵活分片数据并行统一到同一个框架中。
它保留了 Megatron-LM 在 TP/PP 上的成熟优化
（如 TransformerEngine 的 FP8 支持、
交错流水线调度），同时利用 FSDP 的自动参数分片
替代了手写的 ZeRO 实现。
核心特性包括：
\begin{itemize}[noitemsep]
  \item 与 PyTorch DTensor 原生兼容的检查点机制，
        支持跨不同并行配置的检查点迁移。
  \item 针对 GB200 NVL72 系统的优化：
        利用 Blackwell 的 NVLink 5.0（1800\,GB/s 双向）
        和 NVSwitch 4.0 的全互联拓扑。
  \item 对 MoE 架构的一等支持，
        包括专家并行与 FSDP 分片的协同。
\end{itemize}

\textbf{veScale-FSDP} 是字节跳动在 2026 年发布的另一个重要方案，% NEW REF: vescale2026 = veScale-FSDP: Flexible and High-Performance FSDP at Scale, Wang et al., arXiv 2026
解决了标准 FSDP 在\textbf{非逐元素优化器}
（如 Muon、Shampoo）和
\textbf{块级量化训练}（如 FP8/FP4 的 block scaling）
上的兼容性问题。
传统 FSDP 将参数展平为一维张量后分片，
这破坏了矩阵级优化器所需的二维结构信息。
veScale-FSDP 通过结构感知的分片策略保留了参数的原始形状，
使得 Muon 等矩阵级优化器可以与 FSDP 无缝集成：
这对于 Kimi K2、Gemini 等使用 Muon 优化器的前沿模型尤为重要。

\textbf{NVIDIA Megatron Core MoE}（arXiv:2603.07685，2026 年 3 月）% NEW REF: yan2026megatronmoe
专门针对 MoE 架构的训练提出了\textbf{全栈协同设计}方案。
MoE 训练中内存、通信和计算三者之间存在强耦合约束：
增大专家数量会增加内存和通信开销，
减小专家粒度会改变计算特性。
Megatron Core MoE 通过统一的抽象层
将专家并行、张量并行和流水线并行的交互形式化，
为大规模 MoE 训练提供了系统化的工程解决方案。

\subsection{十万卡级别的训练集群}
\label{sec:100k-gpu}

2025--2026 年，训练集群的规模跨入了\textbf{十万卡}量级。

\textbf{xAI Colossus} 是目前最大的单站点 AI 训练集群。
2025 年中期上线时配备了约 100,000 块 NVIDIA H100 GPU，
仅用 122 天就从零建成。
到 2026 年 1 月，Colossus 扩展到 555,000 块 GPU，
总功耗达 2\,GW（相当于两个中型核电站），
硬件投入约 180 亿美元。
这个集群用于训练 Grok 系列模型，
其规模已经\textbf{远超}此前任何已知的训练配置：
Meta 训练 LLaMA 3 405B 使用了 16,384 块 H100，
DeepSeek-V3 仅用了 2,048 块 H800。

在如此规模下运行面临全新的工程挑战：
\begin{itemize}[noitemsep]
  \item \textbf{故障率}：假设单块 GPU 的年故障率为 5\%，
        100K GPU 集群中每天平均有
        $100000 \times 0.05 / 365 \approx 14$ 块 GPU 发生故障。
        系统必须支持\textbf{弹性训练}：
        在不中断整体训练的情况下隔离故障节点并重新分配计算。
  \item \textbf{检查点开销}：100K GPU 的模型状态可能达到数十 TB，
        频繁保存检查点的 I/O 开销不容忽视。
        异步检查点写入和分布式文件系统（如 Lustre、DAOS）
        是必需的基础设施。
  \item \textbf{网络拓扑}：十万卡级别的集群不可能用单一的
        Fat-Tree 拓扑——交换机数量和布线复杂度都不允许。
        xAI 使用了多层 Rail-Optimized + Fat-Tree 混合拓扑，
        并在软件层面根据通信模式动态选择路由路径。
\end{itemize}

\subsection{Blackwell 架构的训练表现}

NVIDIA Blackwell 架构（B200/GB200）在 2025 年下半年
进入规模化部署，MLPerf Training v5.0/v5.1 基准测试
提供了第一批公开的性能数据。

\textbf{GB200 NVL72} 是 Blackwell 的旗舰训练系统：
72 块 B200 GPU 通过第五代 NVLink 全互联
（1.8\,TB/s 双向带宽），
形成一个共享显存池（72 $\times$ 192\,GB $=$ 13.8\,TB）。
在 MLPerf Training v5.1 中，
5,120 块 Blackwell GPU 将 LLaMA 3.1 405B 的训练时间
压缩到\textbf{约 10 分钟}：
相比 Hopper 架构实现了每 GPU $2.2\times$--$2.6\times$ 的加速。

Blackwell 的关键训练加速来自多个维度的\textbf{协同提升}：
不是单一突破，而是“全面翻倍”：

\begin{enumerate}[noitemsep]
  \item \textbf{FP4 硬件原生支持}：
        B200 的第五代 Tensor Core
        原生支持 FP4（NVFP4 格式）矩阵乘法，
        理论 FLOPS 是 FP8 的 $2\times$（9\,PFLOPS FP4 vs 4.5\,PFLOPS FP8）。
        这使得训练中\textbf{计算密集}的操作
        （矩阵乘法占训练总计算的 90\%+）
        可以获得接近 $2\times$ 的加速。
  \item \textbf{NVLink 5.0}：
        双向带宽从 H100 的 900\,GB/s 提升到 1800\,GB/s，
        直接缓解了张量并行的通信瓶颈。
        更重要的是 NVL72 的全互联：
        72 块 GPU 共享显存，
        使得 TP 度数可以远超 8。
  \item \textbf{HBM3e 升级}：
        192\,GB 显存 + 8\,TB/s 带宽，
        使得更大的 batch size 和更长的序列成为可能。
        显存带宽的 $2.4\times$ 提升
        对推理（decode 阶段内存带宽受限）
        的影响甚至大于对训练。
  \item \textbf{Transformer Engine 2.0}：
        Blackwell 的 Transformer Engine
        增加了对 FP4 量化的自动管理，
        包括逐层的精度选择
        （某些对精度敏感的层使用 FP8，
        不敏感的层使用 FP4）
        和动态缩放因子的硬件加速。
  \item \textbf{第二代 Confidential Computing}：
        B200 增加了硬件级别的可信执行环境（TEE），
        允许在\textbf{加密的显存}中训练模型：
        训练数据和模型参数对云服务商不可见。
        这对于需要在第三方云上训练
        但训练数据高度敏感的场景
        （如医疗、金融）具有重要意义。
\end{enumerate}

\subsubsection{Blackwell 的功耗和散热挑战}

Blackwell 的性能提升伴随着\textbf{功耗的显著增加}。
单块 B200 的 TDP（热设计功耗）为 \textbf{1000W}：
远超 H100 的 700W。
一个 GB200 NVL72 机架的总功耗约为 \textbf{120\,kW}：
这在传统数据中心中是\textbf{不可接受的}。
作为参考，传统的企业服务器机架功耗通常为 5--15\,kW，
H100 DGX 机架约为 10.2\,kW。
Blackwell 的功耗密度是传统服务器的 \textbf{10 倍以上}。

这使得\textbf{液冷}（liquid cooling）
从“可选的效率优化”变成了\textbf{必需的基础设施}。
GB200 NVL72 使用\textbf{直接液冷}（Direct-to-Chip Liquid Cooling，D2C）：
冷却液直接流过 GPU 芯片表面的铜板，
而非通过空气间接散热。
D2C 的散热效率是空气冷却的 $3\times$--$5\times$，
且允许更高的功率密度。
但 D2C 需要全新的数据中心设计：
管道、泵组、热交换器、冗余系统：
改造现有数据中心的成本约为
\textbf{每机架 20--40 万美元}。

对于 LLM 训练的经济学而言，
功耗和散热成本不再是可以忽略的“运维细节”：
它们正在成为\textbf{训练总成本}的重要组成部分。
一个 10,000 块 B200 的训练集群，
仅电费一项（按 \$0.08/kWh 计算）
一年就需要约 \textbf{7000 万美元}：
这还不包括散热和设施成本。
这推动了整个行业对\textbf{能效}（performance per watt）
的关注：
不仅要追求绝对的训练速度，
还要优化每瓦特的训练产出。

\subsubsection{DGX vs HGX：两种交付形态}

NVIDIA 的 GPU 训练系统有两种主要交付形态，
理解它们的区别有助于理解训练集群的构建：

\textbf{DGX}（如 DGX H100/B200）是 NVIDIA 的\textbf{完整系统}：
包含 GPU、CPU、内存、存储、NVLink、网卡和软件栈，
开箱即用。DGX 的优势是快速部署和统一的技术支持，
但灵活性有限（无法自定义硬件配置）且价格较高
（DGX H100 售价约 \$300K）。

\textbf{HGX}（如 HGX H100/B200）是 NVIDIA 的\textbf{GPU 基板}：
只包含 GPU 和 NVLink/NVSwitch，
需要系统集成商（如 Dell、HPE、Supermicro）
将其安装到自己设计的服务器中。
HGX 的优势是灵活性（可以选择不同的 CPU、内存、网卡组合）
和更低的价格（HGX H100 约 \$200K，节省约 \$100K/台）。
大多数大规模训练集群使用 HGX 而非 DGX：
因为在数千台机器的规模下，
每台节省 \$100K 意味着数亿美元的总成本差异。

Meta、Google、Microsoft 和 xAI 的超大规模集群
都使用了基于 HGX 的自定义设计。
Meta 的 Grand Teton（用于 LLaMA 3 训练）
基于 HGX H100 构建了自有的 Open Rack 设计，
并开源了服务器的硬件规格：
这是 Open Compute Project (OCP) 的一部分，
目标是通过开放硬件设计来降低 AI 基础设施的成本。

\section{混合精度训练}

\subsection{数值格式详解}

传统的神经网络训练使用 32 位浮点数（FP32）。
但研究发现，大多数训练计算可以用更低精度完成而不损失质量。
理解不同的数值格式是混合精度训练的基础。

IEEE 浮点数的表示为：
$(-1)^s \times 2^{e - \mathrm{bias}} \times (1 + m)$，
其中 $s$ 是符号位，$e$ 是指数，$m$ 是尾数。
不同格式的区别在于分配给指数和尾数的位数：

\begin{center}\small
\renewcommand{\arraystretch}{1.3}
\begin{tabular}{lccccl}
  \toprule
  \rowcolor{figLightGray}
  \textbf{格式} & \textbf{总位数} & \textbf{符号} & \textbf{指数} & \textbf{尾数}
  & \textbf{数值范围} \\
  \midrule
  FP32 & 32 & 1 & 8 & 23 & $\sim 10^{\pm 38}$ \\
  FP16 & 16 & 1 & 5 & 10 & $\sim 6.5 \times 10^4$ \\
  BF16 & 16 & 1 & 8 & 7 & $\sim 10^{\pm 38}$ \\
  FP8 E4M3 & 8 & 1 & 4 & 3 & $\sim 448$ \\
  FP8 E5M2 & 8 & 1 & 5 & 2 & $\sim 57344$ \\
  \bottomrule
\end{tabular}
\end{center}

\textbf{BF16}（Brain Float 16）有 8 位指数和 7 位尾数。
关键洞察：指数位数与 FP32 完全相同，
因此数值范围也完全相同（$\sim 10^{38}$）。
代价是尾数从 23 位减少到 7 位：
精度大约等于 $2^{-7} \approx 0.78\%$ 的相对误差。
但在训练中，随机梯度下降本身就是一个噪声过程，
0.78\% 的精度损失几乎可以忽略。

\textbf{FP16}（半精度）有 5 位指数和 10 位尾数。
它的精度比 BF16 更高（$2^{-10} \approx 0.1\%$ 相对误差），
但数值范围只到 $\sim 65504$。
这在训练中是致命的：loss spike 时梯度可能达到 $10^5$ 量级，
直接溢出为 \texttt{inf}，导致训练崩溃。

\begin{whybox}[为什么 BF16 优于 FP16？]
FP16 的数值范围只到 $\sim 65504$，
在训练中很容易溢出（loss spike 时梯度可能非常大）。
BF16 的范围与 FP32 相同，几乎不会溢出。
虽然 BF16 的精度（7 位尾数 vs.\ 10 位）稍低于 FP16，
但这在训练中影响很小，
因为随机梯度下降本身就是一个「噪声」过程。
BF16 是 Google TPU 原生支持的格式，
NVIDIA 从 A100 开始也全面支持 BF16。
现在所有主流 LLM 的预训练都使用 BF16。
\end{whybox}

\subsection{FP16 的 Loss Scaling}

如果必须使用 FP16 训练
（例如在不支持 BF16 的旧硬件上），
需要 \textbf{Loss Scaling} 来避免梯度下溢。

问题在于：
FP16 能表示的最小正规格化数约为 $6 \times 10^{-5}$，
但训练后期很多梯度的量级在 $10^{-7}$ 到 $10^{-8}$，
这些梯度在 FP16 中会被截断为 0（下溢）。
解决办法是在计算损失时乘以一个大的缩放因子 $S$
（通常 $S = 2^{16}$ 或更大），
使得反向传播中的梯度整体放大 $S$ 倍，
避免下溢。
在更新参数前再除以 $S$ 恢复真实的梯度值。

动态 Loss Scaling（Dynamic Loss Scaling）
自动调整 $S$：
每 $K$ 步，如果没有梯度溢出就增大 $S$；
如果检测到 \texttt{inf} 或 \texttt{NaN}，就减小 $S$ 并跳过这步更新。

BF16 不需要 Loss Scaling，因为它的数值范围与 FP32 相同：
这是 BF16 被广泛采用的另一个重要原因。

\subsection{AMP：自动混合精度}

PyTorch 的 AMP（Automatic Mixed Precision）
是混合精度训练的标准实现。
核心思想是：
\begin{itemize}[noitemsep]
  \item \textbf{前向传播}：在低精度（BF16/FP16）下进行，减少内存和加速计算。
  \item \textbf{主权重}（Master Weights）：在 FP32 下维护一份完整精度的参数副本。
  \item \textbf{梯度更新}：在 FP32 下进行，将低精度梯度累加到 FP32 主权重上。
\end{itemize}

为什么需要 FP32 主权重？
因为参数更新通常很小（$\mathrm{lr} \times \mathrm{grad} \sim 10^{-7}$），
如果直接在 BF16 中更新，
$\mathrm{param} + \delta$ 可能等于 $\mathrm{param}$：
更新被精度「吃掉」了。
FP32 的高精度确保了微小的更新不会丢失。

\subsection{FP8：DeepSeek 的突破}

DeepSeek-V3 是首个完全使用 FP8 精度进行端到端训练的大规模模型~\cite{deepseekv3}。
FP8 有两种变体：E4M3（4 位指数 + 3 位尾数，用于前向）
和 E5M2（5 位指数 + 2 位尾数，用于反向梯度）。

为什么前向用 E4M3、反向用 E5M2？
前向传播中，激活值的数值分布相对集中，
4 位指数的范围（$\sim 448$）通常够用，
多一位尾数（3 vs 2）带来的精度提升更有价值。
反向传播中，梯度的数值范围波动更大，
5 位指数的更大范围（$\sim 57344$）更重要，
牺牲一位尾数是可接受的。

核心挑战在于精度。FP8 的 3 位尾数只能表示 8 个不同的值：
直接用于矩阵乘法会导致严重的累积误差。
DeepSeek 的解决方案是\textbf{细粒度量化}：
将权重矩阵切成 $128 \times 128$ 的小块，
每个小块有独立的缩放因子；
激活值则以 $1 \times 128$ 的粒度量化。

为什么要用 $128 \times 128$ 的块？
因为同一个矩阵中，不同区域的数值分布可能差异很大：
如果用整个矩阵的统一缩放因子，
数值小的区域会被压缩到 FP8 的最低几个值，
丢失大量信息。
块级缩放让每个小块有自己的「标尺」，
大幅提高了量化的保真度。

再加上每 4 次 WGMMA 操作后强制 FP32 累加，
将累积误差控制在可接受范围内。
WGMMA（Warp Group Matrix Multiply-Accumulate）是
NVIDIA Hopper 架构（H100）引入的 Tensor Core 指令：
FP8 的乘法在 Tensor Core 中完成，
但部分和定期累加到 FP32 寄存器中，
防止小数值被大数值「淹没」。

FP8 训练将显存需求减半，吞吐量提升 30--40\%。
PyTorch 原生的 \texttt{torchao.float8} 包
在 2K H200 集群上验证了 1.34--1.43$\times$ 的训练加速。

\begin{warnbox}[FP8 训练的稳定性挑战]
FP8 训练并非无风险。
DeepSeek-V3 的技术报告提到，
FP8 训练在某些条件下会出现更频繁的 loss spike，
尤其是在训练早期模型参数分布还不稳定时。
关键的稳定化措施包括：
\begin{enumerate}[noitemsep]
  \item 训练早期（前几千步）使用 BF16 预热，之后切换到 FP8。
  \item 对归一化层（RMSNorm）和 Attention 的 softmax
        始终保持 FP32/BF16，这些操作对精度极为敏感。
  \item 动态缩放因子的更新需要延迟：
        使用上一步的统计量来设定当前步的缩放因子，
        避免缩放因子的震荡。
\end{enumerate}
\end{warnbox}

\textbf{$\mu$nit Scaling} 是 2025 年提出的一种简化 FP8 训练的方法。% NEW REF: narayan2025munit = $\mu$nit Scaling: Simple and Scalable FP8 LLM Training, Narayan et al., 2025
传统 FP8 训练需要精心调节动态缩放因子（dynamic scaling factors），
而 $\mu$nit Scaling 通过系统化的参数初始化和梯度归一化
\textbf{完全消除了动态缩放的需要}：
FP8 训练变得和 BF16 训练一样简单，
无需额外的超参数调节。
它还支持跨模型宽度的超参数迁移（hyperparameter transfer），
小模型上调好的超参数可以直接用于大模型。

\subsection{FP4：下一个精度前沿}
\label{sec:fp4-training}

FP4 训练是 2025--2026 年低精度训练的最新前沿。
FP4 仅用 4 位表示一个浮点数：
1 位符号 + 2 位指数 + 1 位尾数（E2M1 格式），
仅能表示 $\{0, 0.5, 1, 1.5, 2, 3, 4, 6\}$ 的有限集合
（加上对应的负数）。
直接用如此粗糙的精度训练似乎不可能：
但 2025 年的研究证明了它的可行性。

\textbf{Quartet}（NeurIPS 2025）证明了% NEW REF: quartet2025 = Quartet: Native FP4 Training Can Be Optimal for Large Language Models, Castro et al., NeurIPS 2025
\textbf{原生 FP4 训练}在理论和实验上都可以接近 BF16 的质量。
关键技巧是\textbf{随机舍入}（Stochastic Rounding, SR）：
将 BF16 数值量化到 FP4 时，
不是简单地取最近值（Round to Nearest），
而是以概率正比于距离进行随机选择。
例如，$1.3$ 有 70\% 的概率被量化为 $1$、
30\% 的概率被量化为 $1.5$：
期望值恰好是 $1.3$，量化是\textbf{无偏}的。
这保证了梯度估计的无偏性，是收敛的关键。

\textbf{NVFP4} 是 NVIDIA 为 Blackwell 架构设计的
专有 FP4 格式。% NEW REF: nvfp4_2026 = Pretraining Large Language Models with NVFP4, NVIDIA, arXiv 2026
与标准 E2M1 不同，NVFP4 使用了优化后的编码方案，
在相同的 4 位预算下提供更好的精度分布。
在 MLPerf Training v5.1 的 LLaMA 3.1 405B 基准测试中，
使用 NVFP4 的 Blackwell GPU 实现了
相比 FP8 Hopper 进一步 $1.5\times$--$2\times$ 的训练加速，
且模型质量与 BF16 基线无显著差异。

\textbf{OCP Microscaling（MX）格式}是另一条重要路线。% NEW REF: mx2025 = An Empirical Study of Microscaling Formats for Low-Precision LLM Training, Yang et al., Meta, 2025
Open Compute Project 定义了一组块级微缩放格式：
每 32 个元素共享一个 E8M0 的缩放因子（8 位指数），
单个元素使用 FP8、FP6 或 FP4 表示。
MXFP4 相比 BF16 实现了 $3.76\times$ 的内存/传输压缩，
$2.73\times$ 的运行时开销降低。

Tseng 等人~(2025) 在 AISTATS 上发表了% NEW REF: tseng2025mxfp4 = Training LLMs with MXFP4, Tseng et al., AISTATS 2025
首个接近无损的 MXFP4 训练配方。
他们发现直接对 MXFP4 使用随机舍入会导致
高方差（来自块级异常值），损害收敛。
解决方案是在量化前应用\textbf{随机 Hadamard 变换}：
将权重和激活值旋转到一个均匀分布的空间中，
理论上界了量化方差，显著改善了训练稳定性。

\begin{center}\small
\renewcommand{\arraystretch}{1.3}
\begin{tabular}{lcccl}
  \toprule
  \rowcolor{figLightGray}
  \textbf{精度} & \textbf{位宽} & \textbf{相对 BF16 加速}
  & \textbf{内存节省} & \textbf{成熟度（2026 年）} \\
  \midrule
  BF16 & 16 bit & $1\times$ & $1\times$ & 工业标准 \\
  FP8 (E4M3/E5M2) & 8 bit & $1.3$--$1.5\times$ & $2\times$ & 广泛采用 \\
  MXFP4 & 4 bit & $2$--$3\times$$^\dagger$ & $\sim$4$\times$ & 实验阶段 \\
  NVFP4 (Blackwell) & 4 bit & $2$--$3\times$ & $\sim$4$\times$ & 早期部署 \\
  \bottomrule
\end{tabular}

\small $^\dagger$\,需要原生 MXFP4 硬件支持（Blackwell），
在 H100 上通过 fake-quantization 模拟则无加速。
\end{center}

\section{完整的训练循环}

理解了分布式训练和混合精度之后，
让我们看一个简化但完整的 LLM 预训练循环：

\begin{lstlisting}[caption={Simplified LLM pre-training loop}]
import torch
from torch.distributed import init_process_group
from torch.nn.parallel import DistributedDataParallel as DDP

# Initialize distributed training
init_process_group(backend="nccl")
local_rank = int(os.environ["LOCAL_RANK"])
torch.cuda.set_device(local_rank)

# Create model with BF16 precision
model = LLM(vocab_size=128000, d_model=4096, n_layers=32)
model = model.to(dtype=torch.bfloat16, device=local_rank)
model = DDP(model, device_ids=[local_rank])

# Optimizer: AdamW with weight decay
optimizer = torch.optim.AdamW(
    model.parameters(),
    lr=3e-4,             # Peak learning rate
    betas=(0.9, 0.95),   # Adam betas
    weight_decay=0.1,    # L2 regularization
)

# Learning rate scheduler: warmup + cosine decay
scheduler = CosineAnnealingWithWarmup(
    optimizer,
    warmup_steps=2000,      # Gradually increase LR
    total_steps=1000000,    # Total training steps
    min_lr=3e-5,            # 10% of peak LR
)

# Training loop
for step, batch in enumerate(dataloader):
    tokens = batch["input_ids"].to(local_rank)  # (B, T)

    # Forward pass
    logits = model(tokens[:, :-1])  # Predict next tokens

    # Cross-entropy loss
    loss = F.cross_entropy(
        logits.reshape(-1, logits.size(-1)),  # (B*T, V)
        tokens[:, 1:].reshape(-1),            # (B*T,)
    )

    # Backward pass
    loss.backward()

    # Gradient clipping (prevent exploding gradients)
    torch.nn.utils.clip_grad_norm_(model.parameters(), 1.0)

    # Update weights
    optimizer.step()
    scheduler.step()
    optimizer.zero_grad()

    # Logging and checkpointing
    if step % 100 == 0:
        print(f"Step {step}, Loss: {loss.item():.4f}")
    if step % 5000 == 0:
        save_checkpoint(model, optimizer, step)
\end{lstlisting}

这段代码展示了预训练的核心循环。几个关键细节值得注意：

\textbf{AdamW 优化器}一直是 LLM 预训练的标准选择
（尽管 2025 年出现了强有力的挑战者——详见 \cref{sec:new-optimizers}）。
它结合了 Adam 的自适应学习率
（每个参数有独立的学习率，基于梯度的一阶和二阶动量）
和权重衰减（L2 正则化，防止参数过大）。
$\beta_1 = 0.9$ 和 $\beta_2 = 0.95$ 是经过大量实验确定的标准值：
$\beta_2$ 比 Adam 原始论文的 0.999 更小，
因为 LLM 训练中数据分布不断变化，
需要更快地「遗忘」旧的梯度统计量。

\textbf{AdamW 与 Adam 的区别}值得展开。
原始的 Adam 将权重衰减作为 L2 正则项加到损失函数中：
$\tilde{L} = L + \frac{\lambda}{2}\|w\|^2$，
这导致权重衰减和自适应学习率之间存在耦合：
当某个参数的梯度很小（Adam 会给它更大的有效学习率）时，
L2 正则的效果也会被放大，行为不直观。
AdamW 将权重衰减\textbf{解耦}：
$w_{t+1} = (1 - \lambda \eta) w_t - \eta \hat{m}_t / \sqrt{\hat{v}_t}$，
直接在参数更新步骤中减去 $\lambda \eta w_t$。
这使得权重衰减的行为更加可预测，
在 LLM 训练中几乎所有团队都使用 AdamW 而非 Adam。

\textbf{学习率调度}是训练稳定性的关键。
预热阶段（warmup）从一个很小的学习率（甚至为零）
线性增加到峰值学习率。
这是因为训练初期，模型参数是随机初始化的，
大的学习率可能导致参数剧烈震荡。
预热让模型先「稳住」，找到一个合理的参数区域，
再以正常速率优化。
之后使用余弦退火（cosine annealing）逐渐降低学习率，
帮助模型在训练后期收敛到更平坦的最小值：
平坦的最小值通常泛化性更好。

还有一种日益流行的替代方案：
\textbf{WSD}（Warmup--Stable--Decay）调度。
它将训练分为三个阶段：
短暂的线性 warmup $\to$
长时间的恒定学习率（stable phase）$\to$
最后阶段快速衰减到最小值。
WSD 的优势在于 stable phase 中学习率恒定，
更易于在不同长度的训练中复用超参数；
而且可以在 stable phase 的任意时刻保存检查点，
之后从该检查点开始 decay 阶段：
这为“继续训练”（continual pretraining）提供了极大的灵活性。

\subsubsection{学习率调度的定量对比}

三种主要的学习率调度策略在实践中的差异
值得做一个定量的比较：

\textbf{余弦退火}（Cosine Annealing）：
\begin{equation}
  \eta(t) = \eta_{\min} + \frac{1}{2}(\eta_{\max} - \eta_{\min})
  \left(1 + \cos\left(\frac{\pi \cdot t}{T}\right)\right)
\end{equation}
其中 $\eta_{\max}$ 是峰值学习率，$\eta_{\min}$ 通常取 $\eta_{\max}/10$，
$T$ 是总训练步数。
余弦退火的特点是学习率\textbf{平滑递减}：
训练中期的学习率下降最快，
而训练初期和末期变化较慢。
GPT-3、LLaMA 1/2/3 均使用余弦退火。

\textbf{WSD}（Warmup--Stable--Decay）：
\begin{equation}
  \eta(t) = \begin{cases}
    \frac{t}{T_w} \cdot \eta_{\max} & t < T_w \\
    \eta_{\max} & T_w \leq t < T_s \\
    \eta_{\max} \cdot \left(\frac{T - t}{T - T_s}\right) & t \geq T_s
  \end{cases}
\end{equation}
其中 $T_w$ 是 warmup 结束步数，
$T_s$ 是 stable phase 结束步数。
WSD 的核心优势是\textbf{灵活性}：
stable phase 的长度可以在训练中\textbf{动态决定}。
如果在 stable phase 的第 50\% 发现模型已经足够好，
可以立即开始 decay 阶段；
如果还需要更多训练，可以延长 stable phase。
这种灵活性在预算不确定的场景下非常有价值。
DeepSeek-V3 和 MiniMax-01 等模型采用了 WSD。

\textbf{逆平方根}（Inverse Square Root）：
\begin{equation}
  \eta(t) = \eta_{\max} \cdot \min\left(1, \frac{T_w}{t}\right) \cdot \frac{1}{\sqrt{t / T_w}}
\end{equation}
这是 Transformer 原论文（Vaswani et al., 2017）使用的调度。
学习率在 warmup 后按 $1/\sqrt{t}$ 缓慢衰减。
这种调度在现代 LLM 训练中已很少使用：
余弦退火和 WSD 在几乎所有比较中都优于逆平方根。

在 MiniCPM 的系统性比较中（面壁智能，2024），
WSD 和余弦退火在\textbf{相同总训练步数}下
达到的最终损失几乎相同（差异 $<$0.5\%）。
但 WSD 的优势体现在\textbf{任意步数的中间检查点质量}上：
由于 stable phase 使用恒定学习率，
任何时刻保存的检查点都处于“完全训练”状态
（没有因为余弦曲线中较高的学习率而“过热”）。
这使得 WSD 特别适合需要\textbf{频繁发布中间版本}的场景。

\subsubsection{Warmup 的细节}

Warmup 看起来简单——“从 0 线性增加到峰值”：
但其细节对训练稳定性有显著影响。

\textbf{Warmup 步数}的选择：
LLaMA 3 使用了 2000 步的 warmup，
GPT-3 使用了 375 步，
DeepSeek-V3 使用了 2000 步。
一般规则是：\textbf{模型越大，warmup 越长}：
更大的模型在随机初始化状态下
参数更“不稳定”，需要更多步来“稳住”。
一个经验公式是 warmup 步数 $\approx 0.1\% \sim 0.5\%$ 的总训练步数，
且不少于 200 步。

\textbf{Warmup 的起始学习率}：
多数实现将起始学习率设为 0，
但 PaLM 使用了 $\eta_{\max}/100$ 作为起始值。
后者可以避免在 warmup 的最初几步
因学习率过低而“浪费”训练计算。
实验表明两者对最终模型质量的影响可以忽略：
但 warmup 从 0 开始更稳定。

\textbf{Batch size warmup}：
与学习率 warmup 类似，
某些训练还使用 \textbf{batch size 逐步增大}的策略。
训练初期使用较小的 batch size
（如目标 batch size 的 1/4），
逐步增大到目标值。
小 batch size 在训练初期提供更多的\textbf{梯度噪声}，
有助于模型“探索”参数空间，
找到更好的初始化区域。
GPT-3 从 32K token/batch 线性增大到 3.2M token/batch，
增长了 100 倍。
LLaMA 3 也使用了类似的 batch size warmup。

\textbf{交叉熵损失}中的 \texttt{tokens[:, :-1]} 和
\texttt{tokens[:, 1:]} 体现了因果语言模型的核心：
输入是位置 $1$ 到 $T-1$ 的 token，
目标是位置 $2$ 到 $T$ 的 token。
模型在每个位置预测下一个 token，
损失衡量预测的概率分布与真实 token 之间的差距。

\textbf{梯度累积}（Gradient Accumulation）
是上面代码中没有显式展示但实际训练中不可或缺的技巧。
当 GPU 显存不足以容纳期望的 batch size 时，
可以将一个大 batch 拆成多个 micro-batch，
每个 micro-batch 的梯度累加起来，
达到一定步数后再执行一次 \texttt{optimizer.step()}。
效果等价于用更大的 batch 训练，
但内存只需一个 micro-batch 的量：

\begin{lstlisting}[caption={Gradient accumulation}]
accumulation_steps = 8  # Effective batch = 8 * micro_batch_size

for step, batch in enumerate(dataloader):
    loss = model(batch) / accumulation_steps  # Scale loss
    loss.backward()  # Gradients accumulate in .grad

    if (step + 1) % accumulation_steps == 0:
        torch.nn.utils.clip_grad_norm_(model.parameters(), 1.0)
        optimizer.step()
        optimizer.zero_grad()
\end{lstlisting}

注意损失需要除以 \texttt{accumulation\_steps}：
因为 \texttt{loss.backward()} 会累加梯度，
如果不缩放，累积后的梯度是 $K$ 个 micro-batch 梯度的\textbf{和}
而非\textbf{平均}，等价于用了更大的学习率。

\section{优化器的演进：超越 AdamW}
\label{sec:new-optimizers}

AdamW 统治 LLM 预训练长达 5 年（2019--2024），
但 2025 年出现了真正的挑战者：
\textbf{Muon 优化器}。

\subsection{Muon：矩阵正交化的动量}

Muon（\textbf{Mo}ment\textbf{U}m \textbf{O}rthogonalized
by \textbf{N}ewton-Schulz）% NEW REF: muon2025 = Muon is Scalable for LLM Training, Jordan et al., arXiv 2025
是一种\textbf{矩阵感知}的优化器。
AdamW 对每个参数独立地维护一阶/二阶动量，
将参数矩阵当作展平的向量处理，忽略了矩阵结构。
Muon 则在参数的\textbf{矩阵空间}中工作：

核心步骤是对累积动量矩阵 $M$ 进行
\textbf{Newton-Schulz 正交化}：
\begin{equation}
  M_{\mathrm{orth}} = \mathrm{NS}_k(M)
  \quad \text{s.t.} \quad M_{\mathrm{orth}}^\top M_{\mathrm{orth}} \approx I
\end{equation}
其中 $\mathrm{NS}_k$ 是 $k$ 步 Newton-Schulz 迭代
（通常 $k = 5$--$10$），
将动量矩阵投影到最近的正交矩阵上。
这等价于取动量的\textbf{极分解}中的正交部分：
保留梯度的方向信息，丢弃尺度信息。

为什么正交化有帮助？直觉上，
正交更新确保了参数空间中
各方向被均匀探索，避免了 Adam 在某些方向上
因自适应学习率的缩放而“停滞”的问题。
实验表明，Muon 在相同的训练 FLOP 下
可以达到 AdamW $1.5\times$--$2\times$ 的收敛速度：
即相同质量的模型只需\textbf{一半的训练计算量}。

\textbf{Muon 的工业级采用}在 2025 年迅速扩展：
\begin{itemize}[noitemsep]
  \item Moonshot AI 的 \textbf{Kimi K2}（1T 参数 MoE）
        使用了 Muon 的改进版 \textbf{MuonClip}，
        在 15.5T token 的训练中实现了“零 loss spike”。
  \item PyTorch 2.10 将 Muon 纳入了官方优化器库。
  \item DeepSeek、Meta、OpenAI 等多家实验室
        被报道在内部实验中采用了 Muon 或其变体。
\end{itemize}

\textbf{MuonClip} 是 Kimi K2 团队对 Muon 的工程改进：% NEW REF: kimik2 = Kimi K2: Open Agentic Intelligence, Kimi Team, arXiv 2025
在 Muon 的基础上加入了\textbf{QK-clip 稳定性技巧}：
当 Attention 的 $QK^\top$ 值超过阈值时，
对 $Q$ 和 $K$ 的梯度进行裁剪。
这解决了 Muon 在超大规模训练中偶尔出现的
注意力分数爆炸问题。

\subsection{其他值得关注的优化器}

\textbf{Sophia}（Second-order Clipped Stochastic Optimization）
使用了廉价的 Hessian 对角线估计来设定逐参数的裁剪阈值。
理论上它利用了二阶信息
（参数更新的步长与该参数的损失曲率成反比），
但避免了完整 Hessian 的巨大计算开销。
早期实验声称训练速度是 Adam 的 $2\times$，
但后续的大规模复现研究表明
其优势在大模型上不如预期明显。

\textbf{SOAP}（Spectral-regularized Adam with Projection）
和 \textbf{Shampoo} 属于\textbf{预条件}
（preconditioning）优化器家族，
它们维护参数梯度的 Kronecker 积近似来逼近完整的 Fisher 信息矩阵。
这类方法在理论上更优雅，但计算和内存开销较大。

Wen 等人~(2025) 在斯坦福的一项系统性基准测试中% NEW REF: wen2025optimizers = Fantastic Pretraining Optimizers and Where to Find Them, Wen et al., arXiv 2025
评估了 11 种优化器在标准化的 LLM 预训练场景下的表现。
使用三阶段坐标下降的严格超参数调优后，
他们的关键发现是：
\begin{enumerate}[noitemsep]
  \item \textbf{Muon 是最强的挑战者}：
        在 1B--8B 规模下一致性地优于 AdamW。
  \item \textbf{公平对比至关重要}：
        很多“新优化器优于 Adam” 的结论
        源于对 Adam 超参数的不充分调优。
        当给予 Adam 同等的调优预算后，差距大幅缩小。
  \item Shampoo 和 SOAP 在小规模实验中表现优异，
        但在大规模下的工程成本（额外的矩阵分解操作）
        部分抵消了优化效果。
\end{enumerate}

\begin{keybox}[优化器选择指南（2026 年）]
\begin{center}\small
\renewcommand{\arraystretch}{1.3}
\begin{tabular}{lccl}
  \toprule
  \rowcolor{figLightGray}
  \textbf{优化器} & \textbf{内存/参数} & \textbf{收敛速度}
  & \textbf{适用场景} \\
  \midrule
  AdamW & 12 bytes & $1\times$（基线） & 通用、成熟、风险最低 \\
  Muon & $\sim$8 bytes & $1.5$--$2\times$ & 大规模预训练，需 FSDP 支持 \\
  MuonClip & $\sim$8 bytes & $1.5$--$2\times$ & 超大规模 MoE 训练 \\
  Sophia & 12 bytes & $\sim$1.3$\times$ & 实验性，大规模收益不确定 \\
  \bottomrule
\end{tabular}
\end{center}
\end{keybox}

\section{多 Token 预测}
\label{sec:mtp}

传统的因果语言模型在每个位置只预测\textbf{一个}下一个 token。
\textbf{多 Token 预测}（Multi-Token Prediction, MTP）% NEW REF: mtp2024 = Better \& Faster Large Language Models via Multi-token Prediction, Gloeckle et al., ICML 2024
训练模型同时预测未来的 $K$ 个 token，
通过引入 $K-1$ 个轻量级的预测头来实现。

\subsection{为什么多 Token 预测有用？}

MTP 的动机有三个层面：

\textbf{更丰富的训练信号}。
标准的 next-token 预测只提供“下一个词是什么”的反馈。
MTP 同时要求模型预测第 2、3、...、$K$ 个未来 token，
迫使模型在更深层次上\textbf{规划}：
要准确预测第 3 个 token，
模型必须对未来的语义走向有一个隐式的“计划”。
这种更长视野的训练信号让模型学到了更好的表示。

\textbf{采样效率}。
每个训练样本（序列）提供的有效梯度信号
增加了 $K$ 倍——不增加前向计算量
（因为预测头非常轻量），
却让每个 token 的“教学价值”大幅提升。

\textbf{推理加速}。
训练好的 MTP 头可以直接用于推理阶段的
\textbf{推测解码}（Speculative Decoding）：
额外的预测头生成候选 token，
主模型验证后一次性接受多个 token，
减少了自回归生成的总步数。

\subsection{DeepSeek-V3 的 MTP 实现}

DeepSeek-V3~\cite{deepseekv3} 使用了 $K = 2$ 的 MTP
（即预测未来 2 个 token）。
其实现架构是\textbf{顺序式}的：
每个额外的预测头不仅接收主模型的隐状态，
还接收\textbf{前一个}预测头的输出。
具体而言，第 $k$ 个 MTP 头的输入为：
\begin{equation}
  \mathbf{h}_t^{(k)} = \mathrm{FFN}^{(k)}\bigl(
    \mathrm{Concat}\bigl[\mathbf{h}_t^{(0)},\;
    \mathrm{Embed}(\hat{x}_{t+k-1})\bigr]
  \bigr)
  \label{eq:mtp-head}
\end{equation}
其中 $\mathbf{h}_t^{(0)}$ 是主模型在位置 $t$ 的隐状态，
$\hat{x}_{t+k-1}$ 是前一个头预测的 token 的嵌入。
$\mathrm{FFN}^{(k)}$ 是一个轻量级的前馈网络（通常只有一层）。

训练损失是主预测头和辅助预测头的加权和：
\begin{equation}
  \mathcal{L}_{\mathrm{MTP}} = \mathcal{L}_{\mathrm{NTP}}
  + \sum_{k=1}^{K-1} \lambda_k \cdot \mathcal{L}_{\mathrm{NTP}}^{(k)}
  \label{eq:mtp-loss}
\end{equation}
其中 $\lambda_k$ 是辅助头的损失权重（通常 $\lambda_1 = 0.3$）。
\textbf{推理时丢弃辅助头}：
MTP 是一种纯训练时的技巧，不增加推理的计算成本
（除非用于推测解码，此时辅助头反而\textbf{加速}推理）。

DeepSeek-V3 的实验表明，MTP 带来了
$\sim$0.5\% 的训练损失改善（在 14.8T token 训练结束时），
这听起来很小，但在前沿模型中
每 0.1\% 的损失改善都对应着显著的下游任务性能提升。

\subsection{MTP 训练课程}

Aynetdinov \& Akbik~(2025)% NEW REF: aynetdinov2025mtp = Pre-Training Curriculum for Multi-Token Prediction in Language Models, Aynetdinov \& Akbik, ACL 2025
在 ACL 2025 上提出了 MTP 的\textbf{训练课程}策略：
不是从训练一开始就使用 MTP，
而是先用标准的 NTP 训练一段时间让模型建立基础能力，
之后再引入 MTP 辅助头。
他们发现这种“渐进式”引入在小模型上
可以显著改善 MTP 的效果：
因为训练早期模型的预测非常不准确，
多 token 预测的信号噪声比太低，
反而可能干扰主模型的学习。

具体而言，Aynetdinov \& Akbik 测试了三种引入策略：
(1) 从第一步就启用 MTP（“全程”模式）；
(2) 在训练进行到 20\% 时启用 MTP（“延迟”模式）；
(3) 逐步增加 MTP 的损失权重
$\lambda_k$，从 0 线性增加到目标值（“渐变”模式）。
在 1.4B 参数模型、100B token 的实验中，
“延迟”模式在下游任务上
比“全程”模式提升了约 1.2\%，
“渐变”模式则提升了约 0.8\%。
有趣的是，在 7B 参数模型上，
三种模式的差异大幅缩小到 0.2\% 以内：
这暗示更大的模型对 MTP 引入时机的敏感性更低，
可能是因为大模型在训练早期就能产出
更准确的预测，使得 MTP 信号更早地变得有用。

这一发现对工程实践的启示是：
对于资源有限的小模型训练，
MTP 的引入时机值得仔细调优；
而对于大规模 frontier 模型，
直接从训练开始就启用 MTP
（如 DeepSeek-V3 的做法）是一个安全的默认选择，
无需为引入时机额外花费超参数搜索的预算。

\section{训练稳定性}

训练一个数十亿参数的模型数周到数月，
过程中不出任何问题是不现实的。
训练不稳定通常表现为\textbf{loss spike}：
损失值突然跳升，如果不处理可能导致训练完全崩溃。

\subsection{Loss Spike：原因与应对}

Loss spike 的成因多种多样，常见的包括：

\textbf{数据质量问题}是最常见的原因。
训练数据中混入了乱码、极长的重复文本、
或格式异常的文档，
这些异常数据会产生异常大的梯度，
导致参数在一步之内偏离到不正常的区域。
LLaMA 3~\cite{dubey2024llama3} 的训练日志显示，
他们通过仔细的数据清洗将 loss spike 的频率降低了 90\%。

\textbf{学习率过高}是另一个常见原因，
尤其是在训练初期 warmup 不够充分时。
模型在高学习率下跨过了损失景观中的「悬崖」（loss cliff），
参数跳到了一个完全不同的区域，
损失急剧上升。

\textbf{数值不稳定}也是重要原因。
在低精度训练中（FP16/FP8），
某些中间计算可能超出数值范围。
例如 Attention 计算中 $QK^\top / \sqrt{d}$ 的值可能很大，
softmax 的指数运算进一步放大，导致溢出。
Flash Attention~\cite{dao2022flashattention} 通过在线 softmax
（逐块计算并维护最大值）缓解了这个问题。

\textbf{应对策略}：
\begin{enumerate}[noitemsep]
  \item \textbf{检查点回滚}：发现 loss spike 后，
        回退到最近的正常检查点重新开始。
        这要求频繁保存检查点
        （通常每 500--1000 步保存一次，
        同时保留最近 3--5 个检查点）。
  \item \textbf{跳过异常 batch}：
        检测到梯度范数异常大的 batch 直接跳过，
        不执行参数更新。
  \item \textbf{临时降低学习率}：
        发现 loss 异常上升后，临时将学习率降低 2--10 倍，
        让模型恢复稳定，再逐渐恢复正常学习率。
  \item \textbf{提高数值精度}：
        对关键操作（归一化、softmax）始终使用 FP32。
\end{enumerate}

\subsubsection{Loss Spike 的定量特征}

理解 loss spike 的\textbf{定量特征}
有助于建立有效的自动检测和响应机制。
基于 PaLM、OPT-175B 和 LLaMA 3 等公开训练日志的分析，
loss spike 可以被分为三个严重级别：

\textbf{轻微 spike}（$\Delta L < 0.5$，相对于滑动平均）：
通常由少量异常数据引起，
模型在 50--200 步内自行恢复，
\textbf{不需要人工干预}。
这类 spike 在大规模训练中每天可能发生 1--3 次：
它们是随机梯度下降在高维空间中
“偶尔走到陡峭区域”的正常现象。
梯度裁剪（$\tau = 1.0$）通常足以控制其影响。

\textbf{中度 spike}（$0.5 \leq \Delta L < 2.0$）：
模型在 200--1000 步后\textbf{可能}恢复，
但恢复后的训练轨迹可能偏离最优路径。
建议在检测到中度 spike 后，
临时将学习率降低 5$\times$ 持续 500 步，
然后线性恢复到正常值。
PaLM 的训练中约 60\% 的中度 spike
通过这种“临时降温”策略成功恢复。

\textbf{严重 spike}（$\Delta L \geq 2.0$ 或检测到 NaN）：
模型几乎无法自行恢复，
必须\textbf{回滚到最近的正常检查点}并跳过当前数据。
严重 spike 通常伴随着梯度范数暴增
（超过正常值的 100$\times$ 以上）
和特定层（通常是 Attention 的 QK 乘积或
FFN 的第二个线性层）的参数异常。
回滚后，建议检查并移除引发 spike 的数据 batch：
PaLM 团队发现，\textbf{同一个有问题的 batch
在不同训练阶段都会引发 spike}，
说明问题的根源在数据而非训练动态。

\subsubsection{自动化的训练监控}

2025 年的大规模训练已经不再依赖人工盯守：
\textbf{自动化监控和响应系统}成为标配。
一个典型的监控系统追踪以下指标：

\begin{itemize}[noitemsep]
  \item \textbf{训练损失}：实时损失值及其 100 步滑动平均。
        如果瞬时损失超过滑动平均的 $1.5\times$，触发轻微警报；
        超过 $3\times$，触发严重警报并自动暂停训练。
  \item \textbf{梯度范数}：全局梯度 L2 范数。
        正常训练中梯度范数应在 0.1--10 的范围内
        （取决于模型规模和学习率）。
        范数超过 100 通常预示 spike 即将发生。
  \item \textbf{参数范数}：各层参数的 L2 范数。
        如果某层参数范数突然增大 $2\times$ 以上，
        说明该层的参数被“推到”了异常区域。
  \item \textbf{激活值统计}：
        各层输出的最大值、均值和标准差。
        Attention 层的 $QK^\top$ 最大值
        是\textbf{最敏感的先行指标}：
        它通常在 loss spike 发生前 10--50 步就开始异常上升。
  \item \textbf{学习率和 batch size}：
        确保调度器按计划工作，
        没有因为软件 bug 导致学习率异常。
  \item \textbf{GPU 利用率和温度}：
        硬件异常（如 GPU 过热降频、
        NVLink 故障导致通信速度下降）
        也会间接影响训练稳定性。
\end{itemize}

\textbf{W\&B}（Weights \& Biases）和 \textbf{MLflow}
是最常用的训练监控平台。
LLaMA 3 的训练使用了自定义的监控仪表板，
实时显示上述所有指标，
并在异常时自动向值班工程师发送告警。
DeepSeek 则开发了自有的训练监控系统，
集成了\textbf{自动回滚}功能：
检测到严重异常时自动加载最近的正常检查点
并跳过有问题的数据区段重新开始，
\textbf{无需人工干预}。
这使得 DeepSeek-V3 的 14.8T token 训练
实现了“零人工回滚”的记录。

\subsection{梯度裁剪}

常用的稳定性措施包括：
\textbf{梯度裁剪}将梯度的全局 L2 范数限制在某个阈值内：
\begin{equation}
  \mathbf{g} \leftarrow \begin{cases}
    \mathbf{g} & \|\mathbf{g}\| \leq \tau \\
    \tau \cdot \frac{\mathbf{g}}{\|\mathbf{g}\|} & \|\mathbf{g}\| > \tau
  \end{cases}
  \label{eq:grad-clip}
\end{equation}
其中 $\tau$ 通常设为 1.0。
这意味着每步参数更新的方向保持不变，
但如果梯度「太大」，步长会被缩放回安全范围。
全局范数裁剪（而非逐参数裁剪）的优势在于
保持了不同参数梯度之间的相对比例关系。

\subsection{Z-Loss：稳定 Softmax 的辅助损失}

\textbf{z-loss} 是 PaLM~\cite{chowdhery2022palm} 引入的
一种辅助正则化损失，
专门用于\textbf{稳定 softmax 的输出层}。

问题的根源是：在训练过程中，
模型输出层（最后一个线性层）的 logit 值
可能变得非常大（$|z| > 100$）。
当 softmax 的输入过大时，
其梯度会变得极小（梯度消失），
同时数值精度急剧下降
（大 logit 的 softmax 输出接近 1.0 或 0.0，
梯度中的 $p(1-p)$ 项趋近于 0）。
这会导致输出层的参数“卡住”：
它们几乎收不到有效的梯度更新。

z-loss 通过在损失函数中添加一个
\textbf{惩罚 logit 绝对值}的正则项来解决这个问题：
\begin{equation}
  \mathcal{L}_{\text{z-loss}} = \alpha \cdot \log^2\!\left(
    \sum_{i=1}^{V} \exp(z_i)
  \right)
  \label{eq:z-loss}
\end{equation}
其中 $z_i$ 是 softmax 之前的 logit，
$V$ 是词表大小，$\alpha$ 是一个很小的系数
（PaLM 使用 $\alpha = 10^{-4}$）。

这个损失项的效果是：
如果 logit 的量级过大（$\sum \exp(z_i)$ 很大），
惩罚会增大，梯度会把 logit 压回到合理范围。
由于 $\alpha$ 很小，z-loss 对主训练目标的干扰可以忽略，
但对训练稳定性的改善是\textbf{显著的}：
PaLM 报告称，添加 z-loss 后
loss spike 的频率降低了约 40\%。

z-loss 现在已经成为大规模预训练的\textbf{标准配置}。
DeepSeek-V3、Qwen3 和 LLaMA 4 都使用了 z-loss
或类似的 logit 正则化技巧。
一个有趣的变体是\textbf{pre-softmax layer norm}：
在 softmax 之前对 logit 做 LayerNorm，
直接将 logit 的量级归一化到 $O(1)$。
这种方法的效果与 z-loss 类似，
但实现更简洁（不需要额外的损失项），
正在被越来越多的团队采用。

\subsection{梯度爆炸与梯度消失}

\textbf{梯度爆炸}和\textbf{梯度消失}
是深度神经网络训练的经典挑战。
在动辄 80--126 层的 LLM 中，
这些问题尤为严峻。

\textbf{梯度爆炸}的机制：
反向传播中，梯度在层与层之间\textbf{乘法传递}。
如果每层的 Jacobian 矩阵的最大奇异值 $\sigma_{\max} > 1$，
梯度经过 $L$ 层后会被放大 $\sigma_{\max}^L$ 倍。
对于 $L = 100$，即使 $\sigma_{\max} = 1.05$，
梯度也会被放大 $1.05^{100} \approx 132$ 倍：
足以导致梯度范数爆炸。
在实践中，梯度爆炸通常伴随着 NaN 值的出现，
这会立即终止训练。

\textbf{梯度消失}的机制类似但相反：
如果 $\sigma_{\max} < 1$，梯度经过多层后
被指数级压缩，使得浅层几乎收不到梯度信号。
这在 ReLU 等不饱和激活函数中有所缓解，
但在 LLM 中仍然存在：
特别是在训练初期，当 Attention 权重几乎均匀分布时，
残差连接（residual connection）是防止梯度消失的关键。

现代 LLM 通过以下\textbf{架构和训练技巧}
控制梯度的传播：

\begin{enumerate}[noitemsep]
  \item \textbf{Pre-Norm 架构}：
        将 LayerNorm 放在子层之前（而非之后），
        使得残差路径上的梯度可以\textbf{直接传播}，
        不受归一化层的衰减。
        几乎所有现代 LLM（LLaMA、DeepSeek、Qwen）
        都使用 Pre-Norm。
  \item \textbf{残差连接}：
        $\mathbf{x}_{l+1} = \mathbf{x}_l + f_l(\mathbf{x}_l)$
        确保了梯度有一条“直通”路径，
        不需要经过非线性变换。
        这是深度 Transformer 能够训练的根本原因。
  \item \textbf{梯度裁剪}（详见上文）：
        直接限制梯度范数的上界。
  \item \textbf{参数初始化}：
        GPT-2 引入的“残差缩放”初始化：
        将残差层的输出乘以 $1/\sqrt{2L}$
        （$L$ 是层数），
        确保随着层数增加，
        残差贡献不会无限累积。
        DeepSeek-V3 和 LLaMA 3 都使用了类似的初始化策略。
  \item \textbf{QK-norm}（2025 年趋势）：
        对 Attention 的 $Q$ 和 $K$ 向量做归一化，
        防止 $QK^\top$ 的值过大导致梯度异常。
        Kimi K2 的 MuonClip 就包含了 QK-clip 机制。
\end{enumerate}

\subsection{权重衰减与 AdamW}

权重衰减（Weight Decay）是 LLM 训练中
普遍使用的正则化手段，典型值为 $\lambda = 0.1$。
它的效果是持续地将参数向零「拉回」：
在每步更新中，参数先衰减 $(1 - \lambda \eta)$ 倍，
然后再按梯度方向更新。

直觉上，权重衰减防止了任何单个参数变得「过大」：
如果某个参数对损失的贡献不大（梯度接近零），
权重衰减会把它逐渐压到接近零；
只有对损失有显著贡献的参数才能「抵抗」衰减而保持非零值。
这是一种隐式的特征选择。

值得注意的是，权重衰减通常\textbf{不应用于偏置项和归一化参数}：
这些参数的量级本身就不大，
衰减它们反而会损害模型表达能力。

\subsection{真实案例：大规模训练的挑战}

大规模训练中的不稳定性不是理论上的担忧：
它们在实践中反复出现：

\textbf{PaLM}~\cite{chowdhery2022palm} 在 6144 块 TPU v4 上训练时
经历了 20 多次 loss spike。
Google 的解决方案是：
每次 spike 发生后，回滚到 spike 前约 100 步的检查点，
\textbf{跳过当时的数据 batch}，然后继续训练。
他们发现大多数 spike 与特定数据 batch 相关：
同样的数据 batch 在不同训练阶段都会引发 spike。

\textbf{OPT-175B} 的训练日志~是公开可查的，% NEW REF: zhang2022opt = OPT: Open Pre-trained Transformer Language Models, Zhang et al., arXiv 2022
记录了一段混乱的训练过程：
硬件故障、loss spike、超参数手动调整、
训练中途改变学习率等。
这份真实的训练日志被社区戏称为「OPT 训练日记」，
它生动地展示了大模型训练远非「设置好超参数后等结果」
那么简单：
它更像是一场需要 24/7 值守的「手术」。

\textbf{DeepSeek-V3}~\cite{deepseekv3} 的训练过程
则是一个正面案例：
整个 14.8T token 的训练过程中\textbf{没有发生过}
不可恢复的 loss spike 或需要回滚的情况。
这得益于他们精心设计的训练基础设施
（FP8 细粒度量化、无辅助损失的负载均衡、
DualPipe 调度等）以及严格的数据质量控制。

\section{训练数据调度策略}
\label{sec:data-scheduling}

预训练不是简单地将全部数据打乱后
一遍一遍地喂给模型：
\textbf{数据出现的顺序和频率}
对最终模型质量有显著影响。
数据调度（data scheduling）是训练循环中
一个经常被低估但影响深远的设计维度。

\subsection{数据退火（Data Annealing）}

\textbf{数据退火}是指在训练后期
提高高质量数据的比例：
与学习率退火（降低学习率）同步进行。
这已经成为 2025--2026 年 frontier 模型训练的\textbf{标准实践}。

数据退火的直觉来自人类学习的类比：
想象一个学生准备考试：
前几周广泛阅读各种材料（建立基础知识），
最后几天则集中复习重点内容（强化核心能力）。
LLM 训练的退火策略遵循同样的逻辑。

具体实施中，退火阶段的配置需要精心设计：

\textbf{退火的时机}：太早开始退火（训练的前 50\%）
会导致模型在高质量数据上过拟合，
因为此时模型的基础能力尚不足以
充分“消化”高质量数据的信息；
太晚开始（最后 2--3\%）
则留给退火的训练预算不足，
效果不显著。
实践中的最优窗口是\textbf{训练的最后 5--10\%}。
LLaMA 3 使用了最后 8\%（约 1.2T token），
DeepSeek-V3 在技术报告中提到了
类似比例的退火阶段。

\textbf{退火数据的组成}：
不同类型的高质量数据在退火中的价值不同。
根据 LLaMA 3 和 Qwen3 的消融实验，
退火阶段最有价值的数据类型依次是：
(1) 高质量数学推理数据
（包括详细解题过程的 GSM8K、MATH 变体）；
(2) 代码教程和文档
（如精选的 GitHub README 和技术博客）；
(3) 高质量的百科全书和教科书内容。
有趣的是，\textbf{对话数据}在退火中的价值相对较低：
因为对话能力主要在后训练的 SFT 阶段培养，
预训练退火对其影响不大。

\subsection{重要数据的重放}

\textbf{数据重放}（data replay）
是另一种在训练过程中
对重要数据给予额外关注的策略。

基本思想是：在训练过程中，
对某些特别重要的数据
进行\textbf{多于一次}的训练。
与简单的 multi-epoch 训练不同，
重放策略是\textbf{有选择性的}：
只对特定类型的高价值数据进行重放，
而不是对所有数据均匀重复。

哪些数据值得重放？
一般原则是：\textbf{数据质量越高、在总训练数据中占比越低，
越值得重放}。
具体来说：

\begin{itemize}[noitemsep]
  \item \textbf{高质量数学证明}：
        这类数据在总训练数据中占比通常 $<$5\%，
        但对数学推理能力贡献很大。
        重放 2--3 次的效果显著优于不重放。
  \item \textbf{精选代码文档}：
        如 Python 标准库文档、
        popular library 的 README 和教程。
        这些数据数量有限但质量极高，
        重放可以强化模型对 API 和编程模式的记忆。
  \item \textbf{高质量百科内容}：
        Wikipedia 的特色条目（Featured Articles）
        约占 Wikipedia 总条目数的 0.1\%，
        但在信息密度和写作质量上远超普通条目。
        对这部分数据进行重放
        可以提升模型的\textbf{事实准确性}。
\end{itemize}

\textbf{重放的时机}同样重要。
Muennighoff 等人（2023）的研究表明，
在训练\textbf{后期}重放高质量数据
比在训练初期重放效果更好：
因为后期的模型已经建立了足够的基础能力，
能够从高质量数据中提取更多信息。
这与退火策略的逻辑一致，
实际上可以将数据重放视为退火的一种特殊形式。

\subsection{课程学习在预训练中的应用}

虽然课程学习（从易到难的训练顺序）
的理论吸引力很强，
但在\textbf{大规模预训练}中的实证结果
目前仍然是\textbf{混合的}：
不像在 SFT 阶段那样一致地有效。

2024--2025 年的多项实验揭示了一个有趣的规律：
\textbf{课程学习在小模型上有效，在大模型上效果递减}。
具体来说：
\begin{itemize}[noitemsep]
  \item 在 $\leq$1B 参数的模型上，
        从简单数据到复杂数据的课程
        可以带来 3--5\% 的训练效率提升
        （达到相同性能所需的步数减少 3--5\%）。
  \item 在 7B--13B 参数的模型上，
        效果缩小到 1--2\%。
  \item 在 70B+ 参数的模型上，
        与随机顺序的差异在统计误差范围内。
\end{itemize}

一种解释是：大模型的参数空间足够大，
即使训练数据的顺序是随机的，
模型也有足够的容量来“自适应”地
先学习简单模式、后学习复杂模式。
小模型则因容量有限，
更依赖于外部调度来引导学习顺序。

因此，2026 年的实践共识是：
对于 frontier 级的大规模预训练，
\textbf{数据退火}（最后阶段提高质量）是值得的，
但\textbf{全程课程学习}（整个训练过程中
按难度排序）的工程成本收益比不高。
退火只需要在最后阶段改变采样权重，
实现简单；而课程学习需要预先为每条数据
分配“难度分数”，并设计复杂的调度策略，
工程投入远大于退火但收益相近。

\section{硬件基础设施}

\subsection{GPU 选择与代际跃升}

2025--2026 年的主力训练 GPU 是 NVIDIA H100 和 H200。
这几代 GPU 的性能跃升幅度堪称惊人：

\begin{center}\small
\renewcommand{\arraystretch}{1.3}
\begin{tabular}{lrrrr}
  \toprule
  \rowcolor{figLightGray}
  & \textbf{A100} & \textbf{H100} & \textbf{H200} & \textbf{B200} \\
  \midrule
  显存 & 80\,GB HBM2e & 80\,GB HBM3 & 141\,GB HBM3e & 192\,GB HBM3e \\
  显存带宽 & 2.0\,TB/s & 3.35\,TB/s & 4.8\,TB/s & 8.0\,TB/s \\
  BF16 算力$^\dagger$ & 312\,TFLOPS & 989\,TFLOPS & 989\,TFLOPS & 2.25\,PFLOPS$^\ddagger$ \\
  FP8 算力$^\dagger$ & -- & 1979\,TFLOPS & 1979\,TFLOPS & 4.5\,PFLOPS$^\ddagger$ \\
  NVLink & 600\,GB/s & 900\,GB/s & 900\,GB/s & 1800\,GB/s \\
  架构 & Ampere & Hopper & Hopper & Blackwell \\
  \bottomrule
\end{tabular}

\small
$^\dagger$\,Dense Tensor Core TFLOPS（不含稀疏加速）。
启用结构化稀疏（2:4 sparsity）后理论峰值翻倍，
但训练中通常不使用结构化稀疏。
$^\ddagger$\,B200 数据来自 NVIDIA 官方规格表（SXM 版本），
实际训练吞吐量取决于软件生态成熟度。
\end{center}

从 A100 到 H100，BF16 算力提升了 $\sim$3.2$\times$，
得益于 Hopper 架构引入的第四代 Tensor Core。
H100 还新增了 FP8 支持，进一步翻倍。
H200 保持与 H100 相同的计算芯片，
但将显存升级到 141\,GB HBM3e：
这意味着可以在\textbf{不减少 GPU 数量}的情况下
放下更大的模型或使用更大的 batch size。
B200（Blackwell）则是全面升级：
192\,GB 显存、4.5 PFLOPS FP8 算力、
1800\,GB/s NVLink，甚至支持 FP4 精度。

Google 的 TPU（Tensor Processing Unit）是另一个重要的训练加速器。
TPU v5e 和 v6e 在 Google Cloud 上可用，
Gemini 系列模型就是在 TPU 上训练的。
TPU 的核心优势在于极高的芯片间互联带宽
和与 Google 基础设施的深度集成。

DeepSeek 使用的是 H800——H100 的中国出口限制版本，
NVLink 带宽从 900\,GB/s 降到 400\,GB/s，
跨节点互联也受到限制。
但 DeepSeek 通过出色的软件优化
（如 DeepEP 的通信--计算重叠策略、DualPipe 的流水线调度）
不仅弥补了硬件限制，
甚至在某些效率指标上超越了使用无限制 H100 的西方实验室。
这充分说明了软件工程的力量：
在硬件被限制的情况下，
聪明的算法和精细的系统优化可以带来巨大的效率提升。

\subsection{互联网络：数据中心的「神经系统」}

训练集群的网络拓扑至关重要。
不同层级的通信有不同的带宽和延迟特征：

\textbf{GPU 间通信（节点内）}：
NVLink 是 NVIDIA 的专有互联协议，
H100 提供 900\,GB/s 的双向带宽（每方向 450\,GB/s）。
8 块 H100 通过 NVSwitch 全互联：
任意两块 GPU 之间都有 900\,GB/s 的带宽，
不存在带宽瓶颈。
NVSwitch 充当 GPU 之间的「交换机」，
支持 All-Reduce、All-Gather 等集合通信操作的硬件加速。

\textbf{节点间通信}：
InfiniBand（IB）是高性能计算的标准互联。
NDR（Next Data Rate）标准提供 400\,Gb/s（50\,GB/s），
XDR 提供 800\,Gb/s（100\,GB/s）。
RoCE（RDMA over Converged Ethernet）是以太网上的 RDMA 方案，
成本更低但性能和稳定性略逊于 IB。

\textbf{网络拓扑设计}直接影响训练效率。
理解不同拓扑的权衡是集群设计的核心能力。

\textbf{Fat-Tree}（Clos 网络）拓扑保证了任意两个节点间
等带宽的通信路径，但成本极高
（需要大量高端交换机）。
一个 1024 节点（8192 GPU）的 Fat-Tree 集群
需要约 320 台 64 端口的 InfiniBand 交换机，
仅交换机成本就超过 1 亿美元。
Fat-Tree 的优势是\textbf{通信模式无关}：
无论训练框架使用何种并行策略，
任意 GPU 之间都有全速通信路径。
但这种“奢侈”在实践中很少被充分利用：
因为大多数训练通信发生在\textbf{固定的 GPU 组}内
（同一 TP 组、同一 PP 阶段或同一 DP 组），
而非随机的 GPU 对之间。

\textbf{Rail-Optimized} 拓扑是一种更经济的替代方案：
每个节点的第 $i$ 块 GPU 连接到同一个 rail 交换机，
同一 rail 上的 GPU 之间通信最快。
具体来说，在 DGX H100 集群中：
节点 1 的 GPU 0、节点 2 的 GPU 0、...、节点 $N$ 的 GPU 0
都连接到同一个 Rail 0 交换机。
这种拓扑的成本约为 Fat-Tree 的 \textbf{40--60\%}，
因为只需要 8 个 rail 交换机（对应 DGX 中的 8 块 GPU）
而非完整的 Leaf-Spine 层次。

Rail-Optimized 拓扑与张量并行（节点内，NVLink）+
数据并行（同 rail 跨节点，IB）的组合特别契合。
但它对\textbf{流水线并行和专家并行}不友好：
这两种并行策略的通信模式
不一定沿着 rail 方向进行。
这也是为什么 Meta 在训练 LLaMA 3-405B 时
选择了改良的 Fat-Tree 拓扑而非 Rail-Optimized：
因为他们使用了 4D 并行（TP + CP + PP + DP），
通信模式更加复杂。

\textbf{Dragonfly} 拓扑是 2025 年开始在超大集群中
出现的新选择。
它将节点组织为多个“组”（group），
组内使用全互联，组间使用稀疏连接。
Dragonfly 的带宽介于 Fat-Tree 和 Rail-Optimized 之间，
但其\textbf{延迟}特性更优：
对于延迟敏感的流水线并行通信，
Dragonfly 的组内全互联可以提供比 Fat-Tree 更低的跳数。
据报道，xAI 的 Colossus 集群
（555K GPU）使用了 Dragonfly 变体拓扑。

\subsubsection{NVLink 的演进}

NVLink 是 NVIDIA 的专有高速互联协议，
其演进直接影响了分布式训练的可能性边界：

\begin{center}\small
\renewcommand{\arraystretch}{1.3}
\begin{tabular}{lcrl}
\toprule
\rowcolor{figLightGray}
\textbf{版本} & \textbf{GPU} & \textbf{双向带宽} & \textbf{互联规模} \\
\midrule
NVLink 1.0 & P100 & 160\,GB/s & 2 GPU 对等 \\
NVLink 2.0 & V100 & 300\,GB/s & 8 GPU (DGX-1) \\
NVLink 3.0 & A100 & 600\,GB/s & 8 GPU (DGX A100) \\
NVLink 4.0 & H100 & 900\,GB/s & 8 GPU (DGX H100) \\
NVLink 5.0 & B200 & 1800\,GB/s & 72 GPU (NVL72) \\
\bottomrule
\end{tabular}
\end{center}

从 NVLink 4.0 到 5.0 的飞跃值得特别关注：
带宽翻倍（900 $\to$ 1800\,GB/s），
但更重要的是互联规模从 8 GPU 扩展到 \textbf{72 GPU}。
这意味着张量并行（TP）的度数
可以从 8（DGX H100 的限制）
扩展到 72：
理论上可以在不使用流水线并行的情况下
训练参数量达数万亿的模型。
虽然实际中 TP=72 的效率会很低
（通信占比过高），
但 TP=16 或 TP=36 的配置
比 TP=8 + PP 的组合更高效：
因为避免了流水线气泡。

\textbf{NVSwitch} 是 NVLink 的“交换机”。
NVSwitch 3.0（H100）在一个 DGX 系统内
实现了 8 块 GPU 的\textbf{全互联}：
任意两块 GPU 之间都有 900\,GB/s 的直连带宽，
不存在带宽瓶颈。
NVSwitch 4.0（B200）将全互联扩展到 72 块 GPU，
这是 NVL72 系统“统一显存池”概念的硬件基础。

\subsubsection{InfiniBand vs RoCE}

在\textbf{节点间通信}领域，
InfiniBand 和 RoCE 是两大竞争方案：

\textbf{InfiniBand}（Mellanox/NVIDIA）
是高性能计算的传统标准：
\begin{itemize}[noitemsep]
  \item NDR（2022）：400\,Gb/s（50\,GB/s）
  \item XDR（2024）：800\,Gb/s（100\,GB/s）
  \item GDR（计划 2026）：1600\,Gb/s（200\,GB/s）
\end{itemize}
InfiniBand 的核心优势是\textbf{低延迟}（$\sim$1\,$\mu$s）
和\textbf{硬件级别的拥塞控制}：
即使在大量并发通信流的情况下，
延迟仍然可预测且稳定。
NVIDIA 通过收购 Mellanox（2020 年）
将 InfiniBand 纳入了自己的产品线，
进一步巩固了“GPU + 互联”的垂直整合。

\textbf{RoCE}（RDMA over Converged Ethernet）
是基于以太网的 RDMA 方案：
\begin{itemize}[noitemsep]
  \item RoCEv2 运行在标准的 25/100/200/400\,GbE 以太网上。
  \item 成本约为 InfiniBand 的 \textbf{50--70\%}
        （利用现有的以太网基础设施）。
  \item 但在\textbf{拥塞场景}下延迟不如 InfiniBand 稳定：
        以太网的 PFC（Priority Flow Control）机制
        在大规模集群中可能引发“暂停帧风暴”，
        导致全网性能退化。
\end{itemize}

DeepSeek 使用 RoCE 而非 InfiniBand：
这部分是因为美国出口管制
限制了 InfiniBand 交换机的供应。
但 DeepSeek 通过精心的\textbf{软件优化}
（特别是 DeepEP 的通信-计算重叠和
DualPipe 的细粒度调度）
在 RoCE 网络上实现了
接近 InfiniBand 的\textbf{有效带宽利用率}。
这再次证明了：
在硬件受限时，
软件工程可以弥补很大一部分差距。

2026 年的趋势是\textbf{Ultra Ethernet}：
一个由 AMD、Broadcom、Cisco、Meta 等公司
推动的新一代以太网标准，
专门为 AI 训练优化。
Ultra Ethernet 在标准以太网的基础上
增加了\textbf{有序可靠传输}和
\textbf{硬件级多路径}支持，
目标是在保持以太网成本优势的同时
达到 InfiniBand 的性能水平。
如果 Ultra Ethernet 成功，
它可能在 2027--2028 年
打破 NVIDIA/InfiniBand 在集群互联领域的主导地位。

\subsection{集群设计与成本}

一个典型的 DGX H100 系统包含 8 块 H100 GPU，
售价约 $\$$250K--$\$$300K（约 $\$$30K--$\$$37K per GPU）。

训练一个 frontier LLM 需要多少钱？

Meta 的 LLaMA 3-405B~\cite{dubey2024llama3}
使用 16384 块 H100 GPU 训练了约 54 天，
仅 GPU 租赁成本估计超过 1 亿美元。
而 DeepSeek-V3 仅用 2048 块 H800，
训练 2 个月，成本约 560 万美元：
性能相当，成本降低了 \textbf{约 20 倍}。

这个巨大的成本差距来自 MoE 架构（每个 token 只激活 5.5\% 的参数）、
FP8 训练（减少内存和计算需求）以及高效的工程实现。
它证明了：在 LLM 训练中，
\textbf{聪明的工程设计可以弥补（甚至超越）硬件优势}。

\textbf{训练成本的详细分解}有助于理解
“每美元训练质量”的来源。
以 LLaMA 3-405B 的训练为例，
成本可以分解为以下几个部分：

\begin{center}\small
\renewcommand{\arraystretch}{1.3}
\begin{tabular}{lrl}
\toprule
\rowcolor{figLightGray}
\textbf{成本项} & \textbf{估算金额} & \textbf{占比} \\
\midrule
GPU 计算（16384 H100 $\times$ 54 天） & $\sim$\$50M & 50\% \\
电力（$\sim$10\,MW $\times$ 54 天） & $\sim$\$9M & 9\% \\
人力（团队 $\sim$50 人 $\times$ 6 个月） & $\sim$\$15M & 15\% \\
数据处理（CPU 计算 + 存储） & $\sim$\$5M & 5\% \\
集群维护（冷却、网络、运维） & $\sim$\$8M & 8\% \\
后训练（RLHF + SFT + 评估） & $\sim$\$8M & 8\% \\
其他（软件许可、cloud 备份等） & $\sim$\$5M & 5\% \\
\midrule
\textbf{总计} & $\sim$\textbf{\$100M} & 100\% \\
\bottomrule
\end{tabular}
\end{center}

（上述数字为基于公开信息和行业估计的近似值。）

几个值得注意的点：
(1) GPU 计算成本虽然是最大项，
但“只有” 50\%，其余 50\% 是
电力、人力、基础设施等“非 GPU” 成本。
这意味着仅仅优化 GPU 利用率（提高 MFU）
最多只能影响总成本的一半。
(2) \textbf{人力成本}占 15\% 并不低：
大模型训练需要高度专业化的团队
（分布式系统工程师、GPU 编程专家、
数据科学家、MLOps 工程师等），
这些人才的薪资远高于行业平均水平。
(3) 电力成本虽然“仅” 9\%，
但在能源价格较高的地区（如欧洲）
可能翻倍到 15--20\%。
这也是为什么 xAI 将 Colossus 集群
建在电力成本较低的 Memphis：
仅电费一项就可能节省数千万美元。

\textbf{DeepSeek-V3 的成本奇迹}需要特别分析。
560 万美元的训练成本
是基于 2048 块 H800 $\times$ 2 个月计算的，
但这\textbf{不包括}：
数据处理、前期实验和消融研究
（DeepSeek 的数据处理流水线和架构探索
可能已经花费了数倍于训练本身的计算）、
人力成本（DeepSeek 的工程团队
以其精英水平闻名）、
以及后训练（RLHF、SFT）的成本。
如果将所有成本都计入，
DeepSeek-V3 的“全周期成本”
可能在 2000--3000 万美元：
仍远低于 LLaMA 3-405B 的 1 亿美元，
但差距从 “20 倍” 缩小到 “3--5 倍”。
这个更现实的比较仍然令人印象深刻，
但也提醒我们不要被单一成本数字所误导。

\textbf{云端 vs 自建}的经济学也值得考量。
按云端 H100 的典型租价 $\sim$$\$$2--3/GPU·hour 计算，
16384 块 GPU 训练 54 天的成本约为
$16384 \times 54 \times 24 \times \$2.5 \approx \$53$M。
自建集群的前期投入更大
（16K H100 的硬件成本约 $\$$500M，加上电力、冷却、机房等），
但如果集群能持续运行多个训练任务，
长期来看单次训练的分摊成本更低。
这就是为什么 Meta、Google、DeepSeek 等
频繁训练模型的机构都选择自建：
而初创公司或偶尔训练的团队更倾向于云端租用。

\subsection{弹性训练与检查点管理}

在数千块 GPU 上运行数周到数月，
\textbf{硬件故障是必然事件而非偶然事故}。
弹性训练和检查点管理是保证训练成功完成的关键基础设施。

\subsubsection{故障率的现实}

硬件故障的频率超出大多数人的直觉。
LLaMA 3~\cite{dubey2024llama3} 的技术报告披露了
在 16384 块 H100 上 54 天训练期间的故障统计：
\begin{itemize}[noitemsep]
  \item 总计 419 次自动检测到的故障
        （平均每天 $\sim$8 次）。
  \item 其中 148 次需要人工干预（约 35\%），
        其余自动恢复。
  \item 主要故障类型包括：
        GPU 硬件错误（78 次）、
        网络连接中断（58 次）、
        存储系统异常（42 次）、
        软件/驱动崩溃（35 次）、
        电源/散热问题（12 次）。
  \item 训练的\textbf{有效利用率}
        （实际训练时间 / 总墙钟时间）约为 90\%：
        剩余 10\% 被故障恢复、检查点保存和维护占据。
\end{itemize}

这些数字意味着：
在 10\% 的时间里，
集群的部分或全部 GPU 没有在做有用的训练工作。
对于价值数千万美元的训练，
10\% 的时间浪费意味着数百万美元的损失。
因此，最大化\textbf{有效利用率}
是大规模训练工程中的核心 KPI。

\subsubsection{检查点策略}

检查点（checkpoint）是训练恢复的基础：
它保存了模型参数、优化器状态和数据加载器的位置，
使得训练可以从中断点继续。

\textbf{检查点大小}随模型规模急剧增长。
对于一个 70B 参数的模型：
\begin{itemize}[noitemsep]
  \item 模型参数（BF16）：$70 \times 10^9 \times 2 = 140$\,GB
  \item 优化器状态（FP32，AdamW）：
        $70 \times 10^9 \times 12 = 840$\,GB
  \item 总检查点大小：约 \textbf{980\,GB}
\end{itemize}

将近 1\,TB 的数据写入存储系统
需要的时间取决于存储带宽：
使用 Lustre 并行文件系统（聚合带宽 $\sim$100\,GB/s）
约需 10 秒，
但如果使用 NFS 或普通 SSD（$\sim$1\,GB/s），
则需要约 16 分钟。
这就是为什么大规模训练集群
\textbf{必须}配备高性能并行存储系统。

\textbf{异步检查点}（asynchronous checkpointing）
是减少检查点开销的关键技术：
训练不暂停，而是在后台将模型状态
异步复制到存储系统。
具体实现是：在检查点步骤，
先将模型状态复制到 CPU 内存（速度极快，$<$1 秒），
然后继续训练；
CPU 后台线程将数据从 CPU 内存
异步写入持久化存储。
这使得检查点操作对训练速度的影响
从“分钟级暂停”降低到“亚秒级中断”。
PyTorch 的 \texttt{torch.distributed.checkpoint}
和 NVIDIA 的 DCP（Distributed CheckPoint）
都支持异步检查点。

\textbf{检查点频率}的选择是一个权衡：
更频繁的检查点意味着更少的“回滚距离”
（故障后最多丢失的训练步数），
但也意味着更高的 I/O 开销。
典型的策略是每 \textbf{500--1000 步}保存一次，
同时保留最近 3--5 个检查点
（旧的自动删除以节省存储空间）。
对于 LLaMA 3 的训练（每步约 2.5 秒），
1000 步约为 40 分钟：
这意味着故障后最多丢失 40 分钟的训练进度。

\subsubsection{弹性训练框架}

\textbf{弹性训练}（elastic training）
允许训练在部分节点故障后\textbf{自动恢复}，
而不需要人工干预。
核心能力包括：
\begin{itemize}[noitemsep]
  \item \textbf{故障检测}：自动检测 GPU 故障、
        网络中断和软件崩溃。
  \item \textbf{节点隔离}：将故障节点从训练中移除，
        重新分配其工作负载给其他节点。
  \item \textbf{自动回滚}：加载最近的正常检查点，
        跳过有问题的数据段，恢复训练。
  \item \textbf{热替换}：在不停止训练的情况下
        将备用节点加入集群，替换故障节点。
\end{itemize}

PyTorch Elastic（\texttt{torchrun}）
是最广泛使用的弹性训练框架。
它支持节点的动态加入和退出，
通过 \texttt{c10d.Store} 维护分布式状态一致性。
但 \texttt{torchrun} 的弹性能力相对基础：
它主要处理“启动时”的节点发现，
而非“运行时”的故障恢复。

2025 年出现了更成熟的弹性训练方案。
\textbf{Oasis}（字节跳动）和
\textbf{Bamboo}（UC Berkeley）
提供了运行时故障恢复能力：
检测到 GPU 故障后，
在 $<$30 秒内完成节点替换和训练恢复，
而非传统方案中 5--10 分钟的手动重启。
对于 16K GPU 的集群
（每天平均 8 次故障），
这种快速恢复每天可以节省
$8 \times 5 = 40$ 分钟的空闲时间：
相当于训练有效利用率提升约 3\%。

\subsection{GPU 利用率：最重要的效率指标}

在大规模训练中，
最关键的效率指标是 \textbf{MFU}
（Model FLOPs Utilization，模型浮点运算利用率）：
实际的计算吞吐量占 GPU 理论峰值的百分比。

更精确地定义：
\begin{equation}
  \mathrm{MFU} = \frac{\text{模型每步的理论 FLOP 数} / \text{每步实际耗时}}
                      {\text{GPU 理论峰值 FLOPS} \times \text{GPU 数量}}
  \label{eq:mfu}
\end{equation}

分子是「如果所有计算资源都用于模型计算，
每秒能做多少 FLOP」；
分母是硬件的理论上限。
MFU 衡量的是「有多少算力真正用在了模型计算上」，
其余部分被通信等待、内存访问、流水线气泡等「浪费」了。

理想的 MFU 是 100\%，但实际中远达不到。
原因包括：通信等待时间、内存带宽瓶颈、
流水线气泡、数据加载延迟等。
典型的 MFU 值：
\begin{itemize}[noitemsep]
  \item \textbf{优秀}：50--60\%（如 PaLM 在 6144 TPU 上达到 57\%）
  \item \textbf{良好}：40--50\%（大多数精心优化的训练任务）
  \item \textbf{一般}：30--40\%（较大的通信开销或未优化的配置）
  \item \textbf{差}：$<$30\%（严重的瓶颈或配置问题）
\end{itemize}

LLaMA 3~\cite{dubey2024llama3} 报告的 MFU 为 38--43\%
（16384 块 H100 GPU）。
这个数字看起来不高，但考虑到 16K GPU 的规模
以及 DP+TP+PP 三维并行的通信开销，
已经是非常出色的工程成就。

每提升 1\% 的 MFU，在 2048 块 GPU 的训练中
可以节省约 1\% 的训练时间：
对于两个月的训练周期，这意味着节省了约 15 小时的 GPU 时间，
换算成成本是数万美元。
这就是为什么基础设施工程师对每一个百分点都精打细算。

\section{国产算力上的大模型训练}
\label{sec:domestic-training}

2025--2026 年，中国大模型训练的一个重大突破是
\textbf{摆脱对 NVIDIA GPU 的依赖}。
在美国持续收紧半导体出口管制的背景下，
国产 AI 芯片生态从“可用”逐步进化到“能训练前沿模型”。

\subsection{GLM-5：零 NVIDIA 依赖的前沿模型}

2026 年 2 月 12 日，智谱 AI 发布了 GLM-5——% NEW REF: glm5_2026 = GLM-5 Technical Report, Zhipu AI, 2026
一个 744B 参数的大语言模型，
在 28.5T token 上训练完成。
GLM-5 的历史意义在于：
\textbf{它完全在 100,000 块华为昇腾 910B 处理器上训练，
没有使用任何一块 NVIDIA GPU}。
这是第一个达到前沿水平的“零 NVIDIA” AI 模型。

更早的 2026 年 1 月 16 日，
智谱发布的 \textbf{GLM-Image} 是第一个% NEW REF: glmimage_2026 = GLM-Image: First SOTA Multimodal Model on Domestic Chips, Zhipu AI, 2026
在国产芯片上训练的 SOTA 多模态模型，
使用华为昇腾 Atlas 800T A2 服务器和 MindSpore 框架完成训练，
一度登顶 Hugging Face 全球排行榜第一。

\subsection{华为昇腾 910C}

昇腾 910C 是目前最强的国产 AI 训练芯片，
其核心规格为：
\begin{itemize}[noitemsep]
  \item \textbf{算力}：$\sim$800 TFLOPS（FP16），
        接近 NVIDIA A100（312 TFLOPS）的 2.5 倍，
        但仍落后于 H100（989 TFLOPS）约 20\%。
  \item \textbf{显存}：96\,GB HBM2e，带宽 1800\,GB/s。
  \item \textbf{互联}：HCCS（Huawei Cache Coherence System）
        提供节点内 GPU 间互联。
\end{itemize}

华为 2026 年计划生产约 600,000 块 910C
（相比 2025 年近乎翻倍），
加上其他型号，全年国产 AI 芯片总产量预计达到约 160 万块。
昇腾的后续路线图包括 950PR $\to$ 950DT $\to$ 960 $\to$ 970，
持续缩小与 NVIDIA 旗舰的性能差距。

昇腾 910C 与 NVIDIA H100 的性能差距
需要从多个维度来理解，而不是仅看 TFLOPS 数字。
在\textbf{矩阵运算}方面，910C 的 Da Vinci 3.0 Core
通过 Cube Unit 实现 FP16 矩阵乘法，
其峰值算力与 H100 的 Tensor Core 差距约 20\%。
但在\textbf{显存带宽}上差距更大：
HBM2e 的 1800\,GB/s 与 H100 HBM3 的 3350\,GB/s
相差约 46\%，这在推理的 decode 阶段
（受内存带宽限制）尤为显著。
HCCS 互联在节点内 8 块芯片间提供约 392\,GB/s 的聚合带宽：
虽然足以支持基本的张量并行，
但与 NVLink 4.0 的 900\,GB/s 双向带宽仍有明显差距，
限制了节点内张量并行的效率。

\textbf{软件生态}是昇腾相对于 NVIDIA 的最大短板。
CUDA 经过 17 年的发展（2007 年至今），
积累了庞大的开发者社区、成熟的调试工具
（如 Nsight、cuda-gdb）和深度优化的算子库
（cuBLAS、cuDNN、NCCL、Transformer Engine）。
CANN 生态虽然在近两年加速成长：
MindSpore 框架和 PyTorch 的 Ascend 后端已可用：
但在算子覆盖率（特别是非标准操作）、
性能调优工具和社区支持上
仍需 2--3 年才能达到 CUDA 的成熟度。
这正是 GLM-5 训练成功的意义所在：
它证明了昇腾生态在\textbf{特定架构}（Transformer）
和\textbf{特定框架}（MindSpore）的组合下，
已经足以支撑前沿模型的全流程训练。

\subsection{DeepSeek 的昇腾适配}

DeepSeek 已原生支持昇腾平台，
提供\textbf{一行代码的 CUDA $\to$ CANN 转换}：
开发者只需极少的代码改动即可将模型迁移到昇腾硬件。
在推理场景下，昇腾上的性能约为 H100 的 60\%，
通过手写 CANN 算子可进一步优化。

\textbf{CUDA 到 CANN 的技术迁移}远比“一行代码”听起来复杂。
CANN（Compute Architecture for Neural Networks）是华为的
AI 计算框架，对应 NVIDIA 的 CUDA 生态。
在底层算子层面，CUDA 的 cuBLAS/cuDNN 算子
需要被替换为 CANN 的 AscendCL 算子：
大部分标准操作（GEMM、卷积、归一化）已有成熟映射，
但 DeepSeek 架构中的若干定制化操作
（如 Multi-Latent Attention 的低秩投影、
MoE 的 All-to-All 通信优化 DeepEP、
以及 FP8 细粒度量化的 $128 \times 128$ 分块策略）
在 CANN 上缺乏直接对应的实现，
需要通过 TBE（Tensor Boost Engine）
或 Ascend C 语言手写高性能算子。
这些算子的性能调优往往需要数周到数月的工程投入。

\textbf{训练 vs 推理的性能差距}是迁移中最显著的挑战之一。
推理场景下，昇腾 910B/910C 可以达到 H100 约 60--70\% 的性能，
这个差距主要来自显存带宽
（HBM2e 1800\,GB/s vs H100 HBM3 3350\,GB/s）
和软件栈的优化成熟度。
但在\textbf{训练}场景下，差距可能更大（40--55\%），
原因是训练对\textbf{跨节点通信}的要求远高于推理：
HCCS 互联的带宽和延迟特性
与 NVLink + InfiniBand 的组合仍有明显差距，
特别是在需要高频 All-Reduce（数据并行）
和 All-to-All（专家并行）操作的 MoE 训练中。
此外，训练中的混合精度策略（FP8 训练）
在昇腾上的支持尚不如 H100 成熟：
NVIDIA 的 Transformer Engine 在 FP8 量化、
动态缩放因子管理和 Tensor Core 调度上
有着数年的领先优势。

\textbf{其他芯片生态}也在积极追赶。
AMD 的 MI300X（192\,GB HBM3，
1307 TFLOPS BF16）
通过 ROCm 生态与 PyTorch 的原生集成，
已经被若干实验室用于 LLM 训练：
其显存容量甚至超过 H100，
在需要超大 batch size 或超长序列的场景下有独特优势。
Intel 的 Gaudi 3 芯片
虽然在性能上不及 NVIDIA 旗舰，
但通过与 DeepSpeed 和 Hugging Face 的深度集成，
在成本效率上具有竞争力。
在中国，除华为昇腾外，
燧原 SG2050、百度昆仑 R300、壁仞 BR100 等芯片
也在各自的生态中探索大模型训练支持：
但它们的软件生态成熟度与昇腾相比仍有代差。

据报道，DeepSeek 在 V4 的开发过程中
优先向国产芯片厂商提供预发布版本进行适配测试：
这一举动的战略意义不言而喻。
它意味着 DeepSeek 的模型架构设计
正在从“CUDA 优先”转向“硬件无关”的方向发展，
为中国 AI 生态摆脱对单一硬件供应商的依赖
奠定了关键的软件基础。

\subsection{长周期训练的稳定性挑战}

尽管国产芯片在算力上已接近可用水平，
\textbf{长周期训练的稳定性}仍是最大的工程瓶颈。
DeepSeek 研究员金宇晨指出：
“长训练过程中的稳定性是中国芯片面临的最大挑战。”

这与前文讨论的 loss spike 问题密切相关，
但在国产芯片上问题更加严峻：
芯片级的数值精度差异、互联协议的成熟度、
以及驱动和框架软件栈的稳定性，
都可能在数周的连续训练中放大为难以追溯的训练异常。
GLM-5 在 100,000 块 910B 上的成功训练
正是攻克这一难题的里程碑式成就。

\section{Jensen Huang 的 AI “五层蛋糕”}
\label{sec:five-layers}

2024--2026 年，NVIDIA CEO 黄仁勋在多次公开演讲中
阐述了他对 AI 计算基础设施的\textbf{五层架构}愿景。
这个框架不仅是 NVIDIA 的商业战略，
更是理解当前 AI 训练工业化格局的有力透镜。
让我们逐层剖析，并评估 2026 年各层的现状。

\subsection{第一层：硅层（Silicon Layer）}

最底层是\textbf{芯片}本身：
GPU 架构的设计与制造。

NVIDIA 在这一层建立了近乎垄断的地位。
从 Volta（V100，2017）到 Ampere（A100，2020）
到 Hopper（H100，2022）到 Blackwell（B200，2024），
每一代架构都为 LLM 训练带来了质的飞跃。
关键的架构创新包括：

\begin{itemize}[noitemsep]
  \item \textbf{Tensor Core 的演进}：
        V100 首次引入 Tensor Core（FP16 矩阵乘法加速）；
        A100 增加了 TF32 和 BF16 支持；
        H100 增加了 FP8 支持和 Transformer Engine；
        B200 增加了 FP4 支持。
        每一代 Tensor Core 的
        理论 TFLOPS 约提升 2--3$\times$。
  \item \textbf{HBM（高带宽内存）的升级}：
        从 HBM2（V100）到 HBM2e（A100）
        到 HBM3（H100）到 HBM3e（H200/B200），
        带宽从 900\,GB/s 增长到 8\,TB/s。
        对于 LLM 训练中内存带宽受限的操作
        （如大 batch 的注意力计算），
        HBM 带宽的提升直接转化为训练速度的提升。
  \item \textbf{互联技术}：
        NVLink 从 1.0（300\,GB/s）进化到 5.0（1800\,GB/s），
        NVSwitch 从节点内 8 GPU 互联
        扩展到 GB200 NVL72 的 72 GPU 全互联。
        这使得单个“超级 GPU”
        （72 块 B200 通过 NVLink 5.0 连接）
        拥有 13.8\,TB 的共享显存：
        足以在\textbf{不使用任何模型并行}的情况下
        放下大多数 LLM。
\end{itemize}

\textbf{CUDA 护城河}是 NVIDIA 在这一层最深的壁垒。
CUDA 生态经过 17 年的积累（2007 至今），
包含数千个优化的算子库（cuBLAS、cuDNN、NCCL、
cuSPARSE、Thrust 等），
数百万开发者的知识积累，
以及几乎所有 AI 框架（PyTorch、TensorFlow、JAX）
的原生支持。
即使竞争对手在\textbf{硬件}上追平了 NVIDIA，
在\textbf{软件生态}上仍可能落后 3--5 年。
这也是为什么 AMD MI300X 虽然在某些硬件指标上
超越了 H100（如 192\,GB 显存 vs 80\,GB），
但在实际训练中的采用率仍远低于 H100。

\textbf{定制 ASIC}是这一层的另一个重要趋势。
Google TPU（从 v1 到 v6e）是最成功的定制 AI 芯片，
Gemini 系列模型全部在 TPU 上训练。
TPU 的核心优势是\textbf{芯片间互联}：
TPU v5 pod 可以将 8960 块芯片
通过高速 ICI（Inter-Chip Interconnect）
连接为一个逻辑设备，
聚合带宽远超同规模的 GPU 集群。
AWS Trainium（从 Trainium1 到 Trainium3）
则瞄准了\textbf{成本效率}：
提供低于 H100 约 40\% 的每 TFLOP 成本，
通过 Neuron SDK 与 PyTorch 集成。
华为昇腾 910C 是中国在这一层的主要代表
（详见 \cref{sec:domestic-training}）。

\subsection{第二层：系统层（Systems Layer）}

第二层是将芯片组装为\textbf{训练系统}：
包括 DGX/HGX 服务器、SuperPOD 集群、
网络互联和存储系统。

\textbf{DGX H100} 是 NVIDIA 的旗舰 AI 服务器：
8 块 H100 GPU + 2 块 Intel Xeon CPU +
2\,TB DDR5 内存 + 8 块 ConnectX-7 400\,Gb/s NIC，
通过 NVSwitch 实现 8 GPU 全互联。
\textbf{DGX SuperPOD} 将多台 DGX 通过
InfiniBand 组成大规模集群。
一个 SuperPOD 通常包含 32 台 DGX（256 块 H100），
通过 Leaf-Spine 拓扑实现全双工通信。

\textbf{GB200 NVL72} 代表了系统设计的范式转变：
从“多个独立服务器通过网络连接”
转变为“单个巨型计算节点”。
72 块 B200 GPU 通过 NVLink 5.0 全互联（不经过 InfiniBand），
共享 13.8\,TB 显存，
形成了一个\textbf{统一的计算和内存池}。
对软件而言，这 72 块 GPU
可以被视为一块拥有 13.8\,TB 显存的“超级 GPU”：
极大地简化了分布式训练的编程模型。
但 NVL72 的\textbf{功耗}（约 120\,kW per rack）
和\textbf{散热}需求
对数据中心的基础设施提出了严峻挑战：
传统的空气冷却已经不够，
液冷（直接液冷或浸没式冷却）成为标配。

\subsection{第三层：基础模型层（Foundation Model Layer）}

第三层是\textbf{模型训练基础设施}：
将硬件系统转化为训练好的 AI 模型的“工厂”。

黄仁勋在多次演讲中将 AI 训练比喻为
“token 的工厂”（token factory）：
输入是数据，输出是“数字智能”。
这个比喻的核心洞察是：
\textbf{AI 模型的训练过程正在工业化}，
如同汽车制造从手工作坊到流水线工厂的转变。

在这一层，关键的技术包括：
分布式训练框架（Megatron-LM、DeepSpeed、FSDP：
本章的主要内容），
训练流程管理（检查点、弹性训练、日志监控），
以及数据处理流水线（\cref{ch:data} 的内容）。
这一层的核心指标是\textbf{每美元的训练质量}：
即在固定的硬件预算下，
能够训练出多好的模型。
DeepSeek-V3 用 560 万美元（2048 块 H800）
训练出了与 LLaMA 3-405B（1 亿美元，16384 块 H100）
竞争的模型：
这 \textbf{20 倍}的效率差距
主要来自第三层（MoE 架构、FP8 训练、
DualPipe 调度等）的创新，而非硬件优势。

\subsection{第四层：平台层（Platform Layer）}

第四层是将训练好的模型部署为\textbf{可用的服务}：
包括推理优化、模型部署和 API 服务。

NVIDIA 的\textbf{NIM}（NVIDIA Inference Microservices）
提供了预优化的模型推理容器，
支持 TensorRT-LLM 的推理加速
（包括 KV cache 管理、连续批处理、
推测解码、FP8/INT4 量化等）。
\textbf{Triton Inference Server}
则是模型服务的标准基础设施，
支持多模型部署、动态批处理和负载均衡。

这一层的竞争在 2025--2026 年急剧加热：
vLLM（UC Berkeley 开源）以其 PagedAttention 技术
在社区中获得了广泛采用，
SGLang 则通过 RadixAttention 提供了更高效的推理调度。
开源推理引擎的快速发展
正在削弱 NVIDIA 在这一层的\textbf{平台锁定}：
用户不再必须使用 NIM/TensorRT 就能获得高效推理。

\subsection{第五层：应用层（Application Layer）}

最顶层是 AI 的\textbf{行业应用}：
将模型能力转化为业务价值。

黄仁勋的愿景是：每个行业都需要自己的“AI 工厂”：
金融、医疗、制造、交通等行业
都将建立专属的 AI 训练和推理基础设施，
使用行业特定的数据训练行业特定的模型。
这是 NVIDIA 的\textbf{最终商业目标}：
不仅卖 GPU，而是成为每个行业
AI 基础设施的“操作系统”。

\textbf{2026 年各层的成熟度评估}：
\begin{center}\small
\renewcommand{\arraystretch}{1.3}
\begin{tabular}{lll}
\toprule
\rowcolor{figLightGray}
\textbf{层级} & \textbf{NVIDIA 地位} & \textbf{竞争态势} \\
\midrule
硅层 & 垄断性领先 & TPU/Trainium/昇腾追赶中 \\
系统层 & 强领先 & 云厂商自建系统 \\
基础模型层 & 工具提供者 & 开源框架主导 \\
平台层 & 强但被侵蚀 & vLLM/SGLang 等开源替代 \\
应用层 & 生态构建中 & 行业 know-how 决定 \\
\bottomrule
\end{tabular}
\end{center}

这个五层框架揭示了一个重要趋势：
\textbf{NVIDIA 的护城河从底层向上层递减}。
在硅层和系统层，NVIDIA 的领先地位在短期内几乎不可撼动；
但在平台层和应用层，开源社区和行业参与者
正在构建不依赖 NVIDIA 生态的替代方案。
这也解释了 NVIDIA 为何积极推进 NIM 等平台产品：
它需要在上层建立新的“粘性”来维持整个生态的价值。

\begin{keybox}[国产算力的战略意义]
美国半导体出口管制的有效性正在被侵蚀。
从 2024 年的“勉强可用”到 2026 年的“能训练前沿模型”，
国产芯片生态经历了质的飞跃：
\begin{itemize}[noitemsep]
  \item GLM-5 证明了 744B 参数模型可以完全在国产芯片上训练。
  \item DeepSeek 的原生昇腾支持降低了迁移门槛。
  \item 2026 年 160 万块国产 AI 芯片的产能
        为大规模训练提供了硬件基础。
\end{itemize}
这并不意味着国产芯片已完全追平 NVIDIA：
在单芯片性能、软件生态成熟度和长期训练稳定性上，
差距依然存在。
但“能训练”和“不能训练”之间的鸿沟已经被跨越。
\end{keybox}


%% ============================================================
% NEW REF: muennighoff2023scaling = Scaling Data-Constrained Language Models, Muennighoff et al., NeurIPS 2023
% NEW REF: gpipe = GPipe: Efficient Training of Giant Neural Networks using Pipeline Parallelism, Huang et al., NeurIPS 2019
% NEW REF: zhang2022opt = OPT: Open Pre-trained Transformer Language Models, Zhang et al., arXiv 2022
%% --- NEW REFERENCES (March 2026 update) ---
% NEW REF: sardana2024inference = Beyond Chinchilla-Optimal: Accounting for Inference in Language Model Scaling Laws, Sardana et al., ICML 2024
% NEW REF: bian2025inference = Scaling Inference-Efficient Language Models, Bian et al., arXiv 2025
% NEW REF: qin2025synthllm = Scaling Laws of Synthetic Data for Language Models, Qin et al., arXiv 2025
% NEW REF: held2025relative = Relative Scaling Laws for LLMs, Held et al., arXiv 2025
% NEW REF: zhang2026prescriptive = Prescriptive Scaling Reveals the Evolution of Language Model Capabilities, Zhang et al., arXiv 2026
% NEW REF: megatronfsdp = Megatron-FSDP: Performance-oriented FSDP for Megatron-LM, NVIDIA, 2025
% NEW REF: vescale2026 = veScale-FSDP: Flexible and High-Performance FSDP at Scale, Wang et al., arXiv 2026
% NEW REF: narayan2025munit = mu-nit Scaling: Simple and Scalable FP8 LLM Training, Narayan et al., 2025
% NEW REF: quartet2025 = Quartet: Native FP4 Training Can Be Optimal for Large Language Models, Castro et al., NeurIPS 2025
% NEW REF: nvfp4_2026 = Pretraining Large Language Models with NVFP4, NVIDIA, arXiv 2026
% NEW REF: mx2025 = An Empirical Study of Microscaling Formats for Low-Precision LLM Training, Yang et al., Meta, 2025
% NEW REF: tseng2025mxfp4 = Training LLMs with MXFP4, Tseng et al., AISTATS 2025
% NEW REF: muon2025 = Muon is Scalable for LLM Training, Jordan et al., arXiv 2025
% NEW REF: kimik2 = Kimi K2: Open Agentic Intelligence, Kimi Team, arXiv 2025
% NEW REF: wen2025optimizers = Fantastic Pretraining Optimizers and Where to Find Them, Wen et al., arXiv 2025
% NEW REF: mtp2024 = Better & Faster Large Language Models via Multi-token Prediction, Gloeckle et al., ICML 2024
% NEW REF: aynetdinov2025mtp = Pre-Training Curriculum for Multi-Token Prediction in Language Models, Aynetdinov & Akbik, ACL 2025
% NEW REF: glm5_2026 = GLM-5 Technical Report, Zhipu AI, 2026
% NEW REF: glmimage_2026 = GLM-Image: First SOTA Multimodal Model on Domestic Chips, Zhipu AI, 2026
% NEW REF: yan2026megatronmoe = NVIDIA Megatron Core MoE: Full-Stack System Co-Design for MoE Training, NVIDIA, arXiv:2603.07685, 2026
